\chapter{مجوزها (Permissions)}

سیستم‌عامل‌های ریشه گرفته از سنت یونیکس، از جمله لینوکس، برخلاف سیستم‌های تک‌کاربره نظیر \en{MS-DOS}، بر پایه معماری چندکاربره (\en{Multi-user}) و قابلیت چندوظیفه‌ای (\en{Multitasking}) بنا شده‌اند. این ساختار به چندین کاربر اجازه می‌دهد تا به‌طور هم‌زمان از منابع یک سیستم واحد بهره‌برداری کنند. اگرچه یک رایانه شخصی ممکن است در ظاهر تنها یک مجموعه صفحه‌کلید و نمایشگر داشته باشد، اما این محدودیت مانع از فعالیت هم‌زمان سایرین نیست؛ برای نمونه، کاربران می‌توانند از راه دور از طریق پروتکل \en{SSH} (\en{Secure Shell}) به سیستم متصل شوند و حتی برنامه‌های گرافیکی را اجرا کرده و خروجی را بر روی دستگاه شخصی خود مشاهده کنند.

ماهیت چندکاربره در لینوکس، یک افزونه یا قابلیت ثانویه نیست، بلکه ویژگی بنیادینی است که در عمق طراحی هسته (\en{Kernel}) آن نهفته است. این رویکرد ریشه در بستری دارد که یونیکس در آن متولد شد؛ در دوران پیش از ظهور رایانه‌های شخصی، کامپیوترها تجهیزاتی عظیم، گران‌قیمت و متمرکز بودند. در آن مدل رایانشی، یک رایانه مرکزی قدرتمند به چندین ترمینال در نقاط مختلف یک سازمان (مانند مراکز دانشگاهی) سرویس‌دهی می‌کرد و به‌طور هم‌زمان از تعاملات ده‌ها کاربر پشتیبانی می‌نمود.

برای تحقق عملی این مدل، ایجاد سازوکارهای حفاظتی جهت جداسازی حریم کاربران امری حیاتی بود. در واقع، معماری سیستم نباید اجازه دهد که فعالیت یک کاربر منجر به اختلال در کل سیستم شود، یا فردی بتواند بدون اجازه به داده‌های سایر کاربران دسترسی یافته و در آن‌ها مداخله کند.

در این فصل، به بررسی این مؤلفه کلیدی از امنیت سیستم پرداخته و دستورات راهبردی زیر را تحلیل خواهیم کرد:

\begin{itemize}
  \item \cmd{id}: نمایش شناسه‌ها و اطلاعات هویتی کاربر جاری.
  \item \cmd{chmod}: تغییر سطوح دسترسی (خواندن، نوشتن، اجرا) برای فایل‌ها و دایرکتوری‌ها.
  \item \cmd{umask}: تعیین مجوزهای پیش‌فرض برای فایل‌ها و دایرکتوری‌های تازه‌ساخته‌شده.
  \item \cmd{su}: اجرای پوسته با هویت یک کاربر دیگر.
  \item \cmd{sudo}: اجرای دستورات با سطح دسترسی ابرکاربر (\en{Superuser}).
  \item \cmd{chown}: تغییر مالک (\en{Owner}) قانونی یک فایل یا دایرکتوری.
  \item \cmd{chgrp}: تغییر گروه (\en{Group}) مالکیت یک فایل یا دایرکتوری.
  \item \cmd{groupadd}: تعریف و افزودن گروه‌های کاربری جدید به سیستم.
  \item \cmd{usermod}: ویرایش و تغییر ویژگی‌های حساب کاربری.
  \item \cmd{passwd}: تنظیم یا تغییر رمز عبور کاربر.
\end{itemize}

\section{مالکان، گروه‌ها و سایر کاربران (Owners, Groups and Others)}

در جریان بررسی ساختار سیستم در فصل سوم، احتمالاً هنگام تلاش برای مشاهده محتوای فایلی حساس مانند \file{/etc/shadow} با خطاهای زیر مواجه شده‌اید:

\begin{latin}
\begin{lstlisting}
[me@linuxbox ~]$ file /etc/shadow
/etc/shadow: regular file, no read permission
[me@linuxbox ~]$ less /etc/shadow
/etc/shadow: Permission denied
\end{lstlisting}
\end{latin}

بروز این خطاها ناشی از آن است که شما در سطح یک «کاربر عادی»، فاقد مجوزهای لازم برای دسترسی به این فایل هستید.

در مدل امنیتی یونیکس، هر فایل و دایرکتوری به یک کاربر (\en{User}) تعلق دارد که به عنوان مالک (\en{Owner}) شناخته می‌شود. مالک فایل، قدرتی کامل جهت کنترل سطوح دسترسی به دارایی خود را داراست. افزون بر این، کاربران می‌توانند در قالب گروه‌های کاربری (\en{Groups}) سازمان‌دهی شوند؛ گروه‌هایی که شامل یک یا چند عضو هستند. مالک می‌تواند مجوزهای مشخصی را به تمامی اعضای یک گروه معین اعطا کند. در نهایت، مجموعه‌ای از مجوزها نیز برای «سایر کاربران» در نظر گرفته می‌شود که در متون فنی لینوکس با نام \en{Others} یا گاهی \en{World} شناخته می‌شوند. این ساختار سلسله‌مراتبی و تفکیک‌شده، شالوده و بنیان سیستم مدیریت دسترسی در لینوکس را تشکیل می‌دهد.

به منظور بازبینی اطلاعات هویتی و اعتبارسنجی جایگاه خود در ساختار سیستم، از دستور \cmd{id} استفاده می‌شود. این ابزار، شناسه‌های عددی منحصربه‌فرد و عناوین گروه‌هایی را که کاربر در آن‌ها عضویت دارد، استخراج و ارائه می‌کند.

\begin{latin}
\begin{lstlisting}
[me@linuxbox ~]$ id
uid=500(me) gid=500(me) groups=500(me)
\end{lstlisting}
\end{latin}

تحلیل خروجی این دستور شامل مؤلفه‌های زیر است:

\begin{itemize}
  \item \textbf{شناسه کاربری (\en{UID - User ID}):} یک مقدار عددی یکتا که سیستم برای شناسایی هر کاربر از آن استفاده می‌کند. در نمونه فوق، \cmd{uid=500(me)} بیانگر آن است که کاربر \en{me} با شناسه ۵۰۰ در هسته سیستم شناخته می‌شود.
  \item \textbf{شناسه گروه اصلی (\en{GID - Group ID}):} مقداری عددی که گروه اولیه و اصلی کاربر را تعیین می‌کند. در اینجا \cmd{gid=500(me)} نشان می‌دهد که گروه اصلی این کاربر نیز هم‌نام با او (\en{me}) تعریف شده است.
  \item \textbf{گروه‌ها (\en{Groups}):} فهرستی جامع از تمامی گروه‌های (اصلی و ثانویه) که کاربر در آن‌ها عضویت دارد.
\end{itemize}

تفاوت در نتایج خروجی میان توزیع‌های مختلف، نظیر فدورا (\en{Fedora}) و اوبونتو (\en{Ubuntu})، ناشی از سیاست‌های متمایز آن‌ها در تخصیص محدوده‌های عددی و مدیریت دسترسی‌های پیش‌فرض است:

\begin{latin}
\begin{lstlisting}
[me@linuxbox ~]$ id
uid=1000(me) gid=1000(me) groups=4(adm),20(dialout),24(cdrom),25(floppy),29(audio),30(dip),44(video),46(plugdev),108(lpadmin),114(admin),1000(me)
\end{lstlisting}
\end{latin}

در اوبونتو، تخصیص شناسه کاربری (\en{UID}) به‌طور معمول از عدد ۱۰۰۰ آغاز می‌گردد، در حالی که در فدورا این فرآیند اغلب از ۵۰۰ شروع می‌شود. همچنین، اوبونتو به‌صورت پیش‌فرض کاربران را در گروه‌های سیستمی متعددی مانند \cmd{adm} (برای مدیریت گزارش‌ها) یا \cmd{dialout} (برای دسترسی به پورت‌های سریال) قرار می‌دهد تا تسهیلات لازم جهت مدیریت منابع و سخت‌افزارها فراهم گردد. این تفاوت‌ها بازتاب‌دهنده استراتژی‌های گوناگون توزیع‌ها در تنظیمات اولیه سطوح دسترسی است.

اطلاعات هویتی و سطوح دسترسی در لینوکس، همسو با سایر پیکربندی‌های حیاتی سیستم، در قالب فایل‌های متنی ساده (\en{Plain Text}) ذخیره و نگهداری می‌شوند. تعریف حساب‌های کاربری در فایل \file{/etc/passwd} و مشخصات گروه‌ها در فایل \file{/etc/group} صورت می‌پذیرد. در فرآیند ایجاد یک کاربر یا گروه جدید، این فایل‌ها به‌روزرسانی شده و هم‌زمان تغییرات لازم در فایل \file{/etc/shadow} — که وظیفه نگهداری امن و رمزنگاری‌شده گذرواژه‌ها را بر عهده دارد — اعمال می‌گردد.

در ساختار فایل \file{/etc/passwd} برای هر حساب کاربری، فیلدهای اطلاعاتی راهبردی زیر تعریف شده است: نام کاربری (\en{Login Name})، شناسه کاربری (\en{UID})، شناسه گروه اصلی (\en{GID})، نام واقعی یا توضیحات کاربر (\en{GECOS})، مسیر دایرکتوری خانگی (\en{Home Directory}) و در نهایت، پوسته پیش‌فرض ورود به سیستم (\en{Login Shell}).

بررسی محتوای فایل‌های \file{/etc/passwd} و \file{/etc/group} آشکار می‌سازد که فراتر از حساب‌های کاربری عادی، شناسه‌های دیگری نیز در سیستم وجود دارند؛ از جمله حساب ابرکاربر (\en{Superuser}) یا \cmd{root} با شناسه عددی صفر (\en{UID 0}) و مجموعه‌ای از «کاربران سیستمی». در فصول آتی و هنگام تحلیل فرآیندها (\en{Processes})، مشاهده خواهید کرد که بسیاری از این موجودیت‌های غیرانسانی، نقش‌های فعالی در مدیریت سرویس‌های پس‌زمینه سیستم ایفا می‌کنند.

در حالی که سنت دیرین یونیکس بر تخصیص تمامی کاربران عادی به یک گروه عمومی (مانند \en{users}) استوار بود، رویکرد مدرن در لینوکس بر پایه مدل «گروه خصوصی کاربر» (\en{User Private Group}) بنا شده است. در این روش، برای هر کاربر یک گروه اختصاصی و تک‌عضوی (هم‌نام با کاربر) ایجاد می‌شود. این راهکار، دقت در تخصیص مجوزها را ارتقا داده و مدیریت اشتراک‌گذاری فایل‌ها را در محیط‌های چندکاربره ساده‌تر و امن‌تر می‌سازد.

\section{مجوزهای خواندن، نوشتن و اجرا}

حقوق دسترسی به اشیاء (\en{Objects}) در سیستم فایل (از جمله فایل‌ها و دایرکتوری‌ها) بر پایه سه مجوز اصلی تعریف می‌شود: خواندن (\en{Read})، نوشتن (\en{Write}) و اجرا (\en{Execute}). برای تحلیل دقیق نحوه پیاده‌سازی این مجوزها، خروجی دستور \cmd{ls -l} را بررسی می‌کنیم:

\begin{latin}
\begin{lstlisting}
[me@linuxbox ~]$ > foo.txt
[me@linuxbox ~]$ ls -l foo.txt
-rw-rw-r-- 1 me me 0 2025-03-06 14:52 foo.txt
\end{lstlisting}
\end{latin}

ده نویسه ابتدایی در این خروجی، نشان‌دهنده مشخصه‌های فایل (\en{File Attributes}) هستند. اولین نویسه، تعیین‌کننده نوع فایل (\en{File Type}) است. جدول زیر متداول‌ترین انواع فایل در اکوسیستم لینوکس را تبیین می‌کند:

\begin{table}[htbp]
\centering
\caption{انواع فایل}
\label{tab:file-types}
\begin{tabularx}{\textwidth}{c l X}
\toprule
\textbf{مشخصه} & \textbf{نوع فایل} & \textbf{شرح فنی} \\
\midrule
\cmd{-} & فایل معمولی (\en{Regular File}) & رایج‌ترین نوع فایل در سیستم. محتوای آن می‌تواند شامل متن ساده (\en{Plain Text})، داده‌های باینری یا برنامه‌های اجرایی (\en{Executable}) باشد. \\
\addlinespace
\cmd{d} & دایرکتوری (\en{Directory}) & فایلی ویژه که نقشی بنیادین در سازماندهی ساختار سلسله‌مراتبی سیستم فایل ایفا می‌کند. یک دایرکتوری در حقیقت فهرستی از ارجاعات به آی‌نودهای (\en{inode}) فایل‌ها و دایرکتوری‌های درون خود را در بر دارد و به‌این‌ترتیب، امکان دسته‌بندی و مدیریت منطقی داده‌ها را در سطوح مختلف فراهم می‌آورد. \\
\addlinespace
\cmd{l} & پیوند نمادین (\en{Symbolic Link}) & یک اشاره‌گر (\en{Pointer}) به فایل یا دایرکتوری دیگر (که «هدف» یا \en{Target} نامیده می‌شود). مجوزهای نمایش‌داده‌شده برای این نوع فایل (معمولاً \cmd{rwxrwxrwx}) صرفاً جنبهٔ بصری داشته و مجوزهای واقعی (نافذ)، متعلق به خودِ فایل یا دایرکتوری هدف هستند. در صورت حذف فایل هدف، پیوند به وضعیت «شکسته» (\en{Broken Link}) درمی‌آید. \\
\addlinespace
\cmd{c} & فایل دستگاهی کاراکتری (\en{Character Special File}) & واسطی برای دستگاه‌هایی که داده‌ها را به‌صورت جریان مستمر بایت (\en{Stream}) انتقال می‌دهند. در این نوع دستگاه‌ها، داده به‌صورت کاراکتر به کاراکتر و به‌صورت ترتیبی (\en{Sequential}) پردازش می‌شود. از مصادیق این گروه می‌توان به ترمینال‌ها، مودم‌ها و فایل ویژهٔ \file{/dev/null} اشاره کرد. \\
\addlinespace
\cmd{b} & فایل دستگاهی بلوکی (\en{Block Special File}) & واسطی برای دستگاه‌هایی که داده‌ها را در بلوک‌هایی با حجم ثابت (معمولاً ۵۱۲ بایت یا بیشتر) مدیریت می‌کنند. برخلاف دستگاه‌های کاراکتری، در این نوع دستگاه‌ها امکان دسترسی تصادفی (\en{Random Access}) به داده‌ها وجود دارد. نمونه‌های بارز این گروه، هارددیسک‌ها (\en{HDD/SSD}) و حافظه‌های \en{USB} هستند. \\
\bottomrule
\end{tabularx}
\end{table}

نُه نویسه باقی‌مانده از فیلد مشخصات فایل، که در ادبیات فنی تحت عنوان حالت فایل (\en{File Mode}) شناخته می‌شوند، بیانگر مجوزهای خواندن (\cmd{r})، نوشتن (\cmd{w}) و اجرا (\cmd{x}) برای سه دسته کاربری هستند:

\begin{itemize}
  \item \textbf{کاربر (\en{User}):} مجوزهای اختصاصیِ مالک قانونی فایل.
  \item \textbf{گروه (\en{Group}):} مجوزهای ناظر بر اعضای گروهی که فایل به آن تعلق دارد.
  \item \textbf{دیگران (\en{Others}):} مجوزهای عمومی برای تمامی حساب‌های کاربری فعال در سیستم.
\end{itemize}

این سطوح دسترسی به‌صورت متوالی و در قالب سه بلوک سه‌تایی طبق ساختار زیر نمایش داده می‌شوند:

\begin{latin}
\begin{lstlisting}
User   Group  Other
rwx    rwx    rwx
\end{lstlisting}
\end{latin}

\begin{table}[htbp]
\centering
\caption{تأثیر ویژگی‌های مجوزها بر روی فایل‌ها و دایرکتوری‌ها}
\label{tab:perm-effects}
\begin{tabularx}{\textwidth}{l X X}
\toprule
\textbf{مجوز (\en{Attribute})} & \textbf{تأثیر بر فایل‌ها} & \textbf{تأثیر بر دایرکتوری‌ها} \\
\midrule
خواندن (\cmd{r}) & امکان باز کردن فایل و خواندن محتویات آن را فراهم می‌کند. & اجازه استخراج لیست نام فایل‌های موجود در دایرکتوری را می‌دهد. برای مشاهده جزئیات تکمیلی (مانند حجم و مجوزها)، برخورداری از مجوز اجرا (\cmd{x}) نیز الزامی است. \\
\addlinespace
نوشتن (\cmd{w}) & امکان اصلاح محتوا یا بازنویسی کامل (\en{Truncate}) فایل را فراهم می‌سازد؛ با این حال، حذف یا تغییر نام فایل صرفاً با این مجوز ممکن نیست و تابع مجوزهای دایرکتوریِ والد است. & به کاربر اجازه ایجاد، حذف یا تغییر نام فایل‌های درون دایرکتوری را می‌دهد. این قابلیت تنها در صورتی فعال می‌شود که مجوز اجرا (\cmd{x}) نیز برای آن دایرکتوری صادر شده باشد. \\
\addlinespace
اجرا (\cmd{x}) & اجازه می‌دهد فایل به عنوان یک برنامه یا فرآیند اجرا شود. در خصوص اسکریپت‌های متنی، علاوه بر مجوز اجرا، برخورداری از مجوز خواندن (\cmd{r}) جهت تفسیر کدها توسط پوسته ضروری است. & امکان ورود به دایرکتوری (مثلاً با دستور \cmd{cd}) و دسترسی به فراداده‌ها (\en{Metadata}) نظیر خروجی \cmd{ls -l} را فراهم می‌کند. اجرای عملیات‌هایی چون \cmd{cp}، \cmd{rm} و \cmd{mv} بر روی اشیاء داخلی، مستلزم وجود این مجوز است. \\
\bottomrule
\end{tabularx}
\end{table}

\begin{table}[htbp]
\centering
\caption{الگوهای رایج مجوزها}
\label{tab:common-perms}
\begin{tabularx}{\textwidth}{l X}
\toprule
\textbf{ویژگی‌های فایل} & \textbf{مفهوم فنی و سطح دسترسی} \\
\midrule
\cmd{-rwx------} & فایل معمولی خصوصی (با مجوز اجرا): تنها مالک فایل مجاز به خواندن، ویرایش و اجرای آن است. هیچ کاربر یا گروه دیگری در سیستم کمترین دسترسی به این فایل را ندارد. \\
\addlinespace
\cmd{-rw-------} & فایل معمولی خصوصی (بدون مجوز اجرا): مالک اجازه مطالعه و تغییر محتوا را دارد، اما مجوز اجرا سلب شده است. برای سایر کاربران کاملاً غیرقابل دسترس است. \\
\addlinespace
\cmd{-rw-r--r--} & فایل عمومی خواندنی: مالک حق خواندن و نوشتن دارد. اعضای گروه و سایر کاربران سیستم (\en{Others}) صرفاً مجاز به خواندن (\en{Read}) فایل هستند. این الگو برای فایل‌های پیکربندی عمومی رایج است. \\
\addlinespace
\cmd{-rwxr-xr-x} & برنامه اجرایی عمومی: مالک دسترسی کامل دارد. سایرین (گروه و دیگران) می‌توانند برنامه را اجرا کرده یا محتوای آن را بخوانند، اما حق تغییر در آن را ندارند. این الگو استانداردِ برنامه‌ها و اسکریپت‌های اجرایی است. \\
\addlinespace
\cmd{-rw-rw----} & فایل اشتراکی گروهی: فقط مالک و اعضای گروهِ مرتبط، اجازه خواندن و ویرایش فایل را دارند. سایر کاربران سیستم هیچ‌گونه دسترسی به آن نخواهند داشت. \\
\addlinespace
\cmd{lrwxrwxrwx} & پیوند نمادین (\en{Symbolic Link}): تمامی پیوندهای نمادین به‌صورت پیش‌فرض این مجوزهای ظاهری را نمایش می‌دهند؛ اما سطح دسترسی واقعی تابعِ مجوزهای فایلِ هدف (\en{Target}) است. \\
\addlinespace
\cmd{drwxrwx---} & دایرکتوری اشتراکی گروهی: مالک و اعضای گروه اجازه ورود، مشاهده فهرست و مدیریت فایل‌ها (ایجاد، حذف و تغییر نام) را دارند. برای سایرین مسدود است. \\
\addlinespace
\cmd{drwxr-x---} & دایرکتوری محافظت‌شده: مالک دسترسی کامل دارد. اعضای گروه صرفاً می‌توانند وارد شده و محتوا را مشاهده کنند، اما حق هیچ‌گونه تغییری در ساختار فایل‌های دایرکتوری ندارند. \\
\bottomrule
\end{tabularx}
\end{table}

\section{دستور chmod: تغییر حالت فایل}

فرمان \cmd{chmod} (مخفف \en{Change Mode}) ابزار اختصاصی برای اصلاح «حالت فایل» (\en{File Mode}) یا همان تنظیمات مربوط به سطوح دسترسی اشیاء در سیستم فایل است. این دستور به مدیر سیستم یا مالک فایل اجازه می‌دهد تا مجوزهای خواندن (\en{Read})، نوشتن (\en{Write}) و اجرا (\en{Execute}) را برای سه دسته کاربری یعنی «مالک»، «گروه» و «سایرین» پیکربندی نماید. شایان ذکر است که از منظر امنیتی، حق تغییر این ویژگی‌ها منحصراً در اختیار مالک قانونی فایل یا ابرکاربر (\en{Superuser}) قرار دارد.

در ساختار دستور \cmd{chmod} از دو متدولوژی مجزا برای تعیین تغییرات استفاده می‌شود: بازنمایی عددی در مبنای هشت (\en{Octal Representation}) و بازنمایی نمادین (\en{Symbolic Representation}). در ادامه، ابتدا به تشریح منطق محاسباتی در روش عددی هشت‌تایی می‌پردازیم که به دلیل اختصار و دقت، در میان متخصصان لینوکس رایج‌تر است.

\subsection{بازنمایی هشت‌تایی (Octal Representation)}

سیستم عددی هشت‌تایی (\en{Octal}) یا مبنای ۸ و همتای آن، شانزده‌تایی (\en{Hexadecimal}) یا مبنای ۱۶، از الگوهای رایج نمایش داده‌ها در علوم رایانه محسوب می‌شوند. این سیستم‌ها در واقع راهکاری برای غلبه بر پیچیدگی‌های کار با سیستم دودویی (\en{Binary}) هستند. ما انسان‌ها به‌طور پیش‌فرض از سیستم ده‌دهی (\en{Decimal}) یا مبنای ۱۰ استفاده می‌کنیم؛ انتخابی که ریشه در فیزیولوژی بدن و شمارش با ده انگشت دارد. در مقابل، رایانه‌ها به دلیل ماهیت مدارهای الکترونیکی، صرفاً دو حالت «روشن» و «خاموش» (صفر و یک) را درک می‌کنند و بر همین اساس از سیستم دودویی یا مبنای ۲ بهره می‌برند.

در سیستم دودویی، شمارش به صورت زیر انجام می‌شود:

\begin{latin}
\begin{lstlisting}
0, 1, 10, 11, 100, 101, 110, 111, 1000, 1001, 1010, 1011...
\end{lstlisting}
\end{latin}

سیستم هشت‌تایی، که مبنای آن عدد ۸ است، از ارقام 0 تا 7 برای شمارش بهره می‌برد:

\begin{latin}
\begin{lstlisting}
0, 1, 2, 3, 4, 5, 6, 7, 10, 11, 12, 13, 14, 15, 16, 17, 20, 21...
\end{lstlisting}
\end{latin}

در سیستم شانزده‌تایی، که مبنای آن ۱۶ است، از ارقام 0 تا 9 و سپس حروف A تا F برای نمایش مقادیر ۱۰ تا ۱۵ استفاده می‌شود:

\begin{latin}
\begin{lstlisting}
0, 1, 2, 3, 4, 5, 6, 7, 8, 9, A, B, C, D, E, F, 10, 11, 12, 13...
\end{lstlisting}
\end{latin}

هر رقم در سیستم شانزده‌تایی (\en{Hex}) دقیقاً معادل ۴ بیت (یک نیبل) و هر رقم در سیستم هشت‌تایی (\en{Octal}) دقیقاً معادل ۳ بیت از یک عدد دودویی است. دلیل این هم‌ترازی ریاضی به تعداد ارقام در هر مبنای عددی بازمی‌گردد. سیستم شانزده‌تایی (مبنای ۱۶) از ۱۶ رقم (0 تا 9 و A تا F) استفاده می‌کند. برای نمایش ۱۶ حالت مجزا در سیستم دودویی، به ۴ بیت نیاز داریم، زیرا ۲ به توان ۴ برابر با ۱۶ است ($16=4^2$). به همین دلیل، هر رقم در مبنای شانزده دقیقاً معادل یک مجموعهٔ ۴ بیتی (یک نیبل) است. در مقابل، سیستم هشت‌تایی (مبنای ۸) از ۸ رقم (0 تا 7) استفاده می‌کند. برای نمایش ۸ حالت مجزا در مبنای دو، به ۳ بیت نیاز داریم، زیرا ۲ به توان ۳ برابر با ۸ است ($8=3^2$). به همین دلیل، هر رقم در مبنای هشت دقیقاً معادل یک مجموعهٔ ۳ بیتی است.

اگرچه منطق دودویی زبان ذاتی ماشین است، اما تعامل مستقیم با آن برای انسان بسیار دشوار و مستعد خطاست. سیستم‌های هشت‌تایی و شانزده‌تایی به‌عنوان واسطی میان انسان و ماشین عمل می‌کنند تا کار با داده‌های سطح پایین و آدرس‌های حافظه تسهیل شود.

برای درک بهتر این مفهوم، یک مثال کاربردی در زمینه بازنمایی رنگ‌ها را در نظر بگیرید. در اکثر سیستم‌های رایانه‌ای، هر رنگ از ترکیب سه مؤلفه اصلی قرمز (\en{Red})، سبز (\en{Green}) و آبی (\en{Blue}) تشکیل می‌شود که برای هر مؤلفه ۸ بیت اختصاص می‌یابد. این ساختار منجر به ایجاد یک عدد ۲۴ بیتی برای هر پیکسل رنگی می‌شود. به عنوان نمونه، یک طیف خاص از رنگ آبی ممکن است در سیستم دودویی به صورت زیر ذخیره شود:

\begin{latin}
\begin{lstlisting}
010000110110111111001101
\end{lstlisting}
\end{latin}

تعامل با چنین توالی‌های طولانی از صفر و یک برای انسان عملاً غیرممکن و به‌شدت مستعد خطا است. در اینجاست که سیستم‌های عددی جایگزین، کارایی خود را نشان می‌دهند.

همانطور که پیش‌تر اشاره شد، هر رقم در سیستم شانزده‌تایی (\en{Hex}) نماینده ۴ بیت و هر رقم در سیستم هشت‌تایی (\en{Octal}) نماینده ۳ بیت از یک عدد دودویی است. این هم‌ترازیِ ریاضی، امکان فشرده‌سازی و نمایش اعداد بزرگ را به شکلی مختصر و خوانا فراهم می‌کند. بنابراین، همان عدد ۲۴ بیتی مربوط به رنگ آبی می‌تواند به سادگی به یک کد شش‌رقمی در مبنای شانزده تبدیل شود:

\begin{latin}
\begin{lstlisting}
436FCD
\end{lstlisting}
\end{latin}

در این قالب، هر جفت رقم \en{hex} نشان‌دهنده یک مؤلفه رنگی است: \cmd{43} برای قرمز، \cmd{6F} برای سبز و \cmd{CD} برای آبی. این تبدیل، فرآیند مدیریت داده‌های بیتی را برای متخصصان بسیار ساده‌تر و دقیق‌تر می‌کند.

اگرچه امروزه سیستم شانزده‌تایی در دنیای رایانش رواج بیشتری دارد، اما ویژگی منحصربه‌فرد سیستم هشت‌تایی در پوشش دقیق ۳ بیت دودویی، آن را به ابزاری فوق‌العاده کارآمد برای مدیریت سیستم مجوزهای لینوکس تبدیل کرده است؛ چرا که مجوزهای لینوکس دقیقاً در گروه‌های سه بیتی (\cmd{r}، \cmd{w}، \cmd{x}) سازمان‌دهی می‌شوند.

با بهره‌گیری از سیستم هشت‌تایی (\en{Octal})، می‌توان تمامی مجوزهای یک فایل را در قالب یک عدد سه‌رقمی واحد پیکربندی کرد. این متدولوژی انطباق کاملی با ساختار امنیتی لینوکس دارد؛ ساختاری که از سه دسته مجوز (برای مالک، گروه و سایرین) تشکیل شده و هر دسته، سه نوع حق دسترسی شامل خواندن، نوشتن و اجرا را در بر می‌گیرد.

از آنجا که هر رقم در مبنای هشت، دقیقاً معادل یک توالی سه‌بیتی در سیستم دودویی (\en{Binary}) است، می‌توان هر ترکیبی از مجوزهای \cmd{r}، \cmd{w} و \cmd{x} را با یک رقم منحصربه‌فرد نمایش داد. در این منطق محاسباتی، هر مجوز دارای ارزش وزنی مشخصی است: برای «خواندن» عدد 4، برای «نوشتن» عدد 2 و برای «اجرا» عدد 1 در نظر گرفته می‌شود. با جمع بستن این مقادیر، سطح دسترسی نهایی برای هر گروه تعیین می‌گردد.

به عنوان مثال، ترکیب عددی 755 به این معناست که مالک فایل دسترسی کامل ($7=1+2+4$)، و اعضای گروه به همراه سایر کاربران، صرفاً مجوز خواندن و اجرا ($5=1+0+4$) را دارا هستند. این روش به دلیل اختصار و دقت بالا، استاندارد طلایی در مدیریت سرورهای لینوکسی محسوب می‌شود.

\begin{table}[htbp]
\centering
\caption{نگاشت مقادیر هشت‌تایی به حالت‌های فایل}
\label{tab:octal-mapping}
\begin{tabularx}{\textwidth}{c c l X}
\toprule
\textbf{عدد هشت‌تایی} & \textbf{معادل دودویی} & \textbf{حالت فایل (\en{Permissions})} & \textbf{شرح عملیاتی} \\
\midrule
\cmd{0} & \cmd{000} & \cmd{---} & فاقد هرگونه مجوز دسترسی \\
\addlinespace
\cmd{1} & \cmd{001} & \cmd{--x} & منحصراً مجوز اجرا \\
\addlinespace
\cmd{2} & \cmd{010} & \cmd{-w-} & منحصراً مجوز نوشتن \\
\addlinespace
\cmd{3} & \cmd{011} & \cmd{-wx} & مجوز نوشتن و اجرا \\
\addlinespace
\cmd{4} & \cmd{100} & \cmd{r--} & منحصراً مجوز خواندن \\
\addlinespace
\cmd{5} & \cmd{101} & \cmd{r-x} & مجوز خواندن و اجرا \\
\addlinespace
\cmd{6} & \cmd{110} & \cmd{rw-} & مجوز خواندن و نوشتن \\
\addlinespace
\cmd{7} & \cmd{111} & \cmd{rwx} & دسترسی کامل (خواندن، نوشتن و اجرا) \\
\bottomrule
\end{tabularx}
\end{table}

با بهره‌گیری از یک عدد سه‌رقمی در مبنای هشت (\en{Octal})، می‌توان به شکلی فشرده و دقیق، سطوح دسترسی یک فایل را برای مالک (\en{Owner})، گروه (\en{Group}) و سایرین (\en{Others}) پیکربندی کرد. در این ساختار، هر رقم به‌ترتیب، نشان‌دهندهٔ سطح دسترسی یکی از این سه دسته کاربری است.

در سناریوی زیر، نحوه تغییر حالت فایل (\en{File Mode}) با استفاده از این متدولوژی نمایش داده شده است:

\begin{latin}
\begin{lstlisting}
[me@linuxbox ~]$ > foo.txt
[me@linuxbox ~]$ ls -l foo.txt
-rw-rw-r-- 1 me me 0 2025-03-06 14:52 foo.txt
[me@linuxbox ~]$ chmod 600 foo.txt
[me@linuxbox ~]$ ls -l foo.txt
-rw------- 1 me me 0 2025-03-06 14:52 foo.txt
\end{lstlisting}
\end{latin}

در مثال فوق، با اجرای فرمان \cmd{chmod 600 foo.txt} و تخصیص کد عددی 600، دسترسی‌ها بازنویسی شدند. طبق منطق هشت‌تایی، عدد 6 (حاصل جمع $0+2+4$) برای مالک به معنای حق خواندن و نوشتن (\cmd{-rw}) است و دو رقم 0 بعدی، تمامی دسترسی‌های گروه و سایر کاربران را به طور کامل سلب کرده است (\cmd{---}). این پیکربندی یکی از ایمن‌ترین حالات برای فایل‌های حاوی داده‌های حساس شخصی است.

اگرچه درک منطق ریاضی میان اعداد هشت‌تایی و دودویی در ابتدای امر ممکن است دشوار به نظر برسد، اما در سناریوهای عملیاتی، تسلط بر پنج عدد کلیدی زیر برای مدیریت اکثر دسترسی‌ها کفایت می‌کند:

\begin{itemize}
  \item \cmd{7} (معادل \cmd{rwx}): دسترسی کامل (خواندن، نوشتن و اجرا).
  \item \cmd{6} (معادل \cmd{rw-}): دسترسی استاندارد فایل‌های متنی (خواندن و نوشتن).
  \item \cmd{5} (معادل \cmd{r-x}): دسترسی استاندارد دایرکتوری‌ها و اسکریپت‌ها (خواندن و اجرا).
  \item \cmd{4} (معادل \cmd{r--}): دسترسی محدود (فقط خواندن).
  \item \cmd{0} (معادل \cmd{---}): عدم دسترسی مطلق (امنیت کامل).
\end{itemize}

\begin{latin}
\begin{lstlisting}
7 (rwx), 6 (rw-), 5 (r-x), 4 (r--), and 0 (---)
\end{lstlisting}
\end{latin}

فرمان \cmd{chmod} علاوه بر روش عددی، از روش بازنمایی نمادین (\en{Symbolic}) نیز برای اصلاح حالت فایل پشتیبانی می‌کند. این رویکرد به شما اجازه می‌دهد تا تغییرات مورد نظر خود را با عباراتی گویاتر و بصری‌تر بیان کنید. ساختار یک عبارت نمادین از سه مؤلفه اصلی تشکیل شده است:

\begin{itemize}
  \item \textbf{مخاطبان (\en{Who}):} تعیین رده‌ای از کاربران که تغییرات باید بر آن‌ها اعمال شود.
  \item \textbf{عملگر (\en{Action}):} مشخص کردن نوع عملیات (افزودن، حذف یا انتساب مستقیم مجوز).
  \item \textbf{نوع مجوز (\en{Permission}):} تعیین ماهیت دسترسی (خواندن، نوشتن یا اجرا).
\end{itemize}

برای شناسایی رده‌های مختلف کاربران در این سیستم، از حروف اختصاری \cmd{u}، \cmd{g}، \cmd{o} و \cmd{a} استفاده می‌شود که جزئیات آن‌ها در جدول زیر ارائه شده است:

\begin{table}[htbp]
\centering
\caption{نمادهای مخاطب در بازنمایی نمادین}
\label{tab:symbolic-who}
\begin{tabularx}{\textwidth}{c l X}
\toprule
\textbf{نماد} & \textbf{مفهوم فنی} & \textbf{شرح کاربرد} \\
\midrule
\cmd{u} & \en{User} & اشاره به مالک قانونی (\en{Owner}) فایل یا دایرکتوری. \\
\addlinespace
\cmd{g} & \en{Group} & اشاره به گروه مالک (\en{Owning Group}) فایل یا دایرکتوری. \\
\addlinespace
\cmd{o} & \en{Others} & اشاره به سایر کاربرانی که نه مالک هستند و نه در گروه مالک عضویت دارند. \\
\addlinespace
\cmd{a} & \en{All} & مخفف «همه»؛ ترکیبی از سه مورد فوق که تغییرات را به‌طور هم‌زمان بر تمامی رده‌ها اعمال می‌کند. \\
\bottomrule
\end{tabularx}
\end{table}

در صورتی که هیچ نماد مشخصی برای تعیین رده کاربران (مانند \cmd{u}، \cmd{g} یا \cmd{o}) در فرمان لحاظ نشود، تغییرات به‌طور خودکار بر تمامی سطوح اعمال خواهد شد.

عملیات مورد نظر برای اصلاح سطوح دسترسی با استفاده از سه عملگر منطقی زیر انجام می‌شود:

\begin{itemize}
  \item \cmd{+} (\textbf{افزودن}): الصاق یک مجوز خاص به دسترسی‌های فعلی، بدون تغییر در سایر مجوزها.
  \item \cmd{-} (\textbf{حذف}): ابطال یک مجوز خاص از فهرست دسترسی‌های موجود.
  \item \cmd{=} (\textbf{انتساب/تنظیم}): بازنشانی کامل مجوزها به مقدار تعیین‌شده. این عملگر تمامی مجوزهای پیشین را نادیده گرفته و ترکیب جدید را جایگزین می‌کند.
\end{itemize}

\begin{table}[htbp]
\centering
\caption{مثال‌های کاربردی از ترکیب‌های نمادین}
\label{tab:symbolic-examples}
\begin{tabularx}{\textwidth}{l l X}
\toprule
\textbf{عبارت نمادین} & \textbf{مفهوم عملیاتی} & \textbf{شرح عملکرد} \\
\midrule
\cmd{u+x} & افزودن مجوز اجرا به مالک & به مالک فایل حق اجرای آن را اعطا می‌کند. \\
\addlinespace
\cmd{u-x} & حذف مجوز اجرا از مالک & حق اجرای فایل را از مالک سلب می‌کند. \\
\addlinespace
\cmd{+x} & افزودن مجوز اجرا به همگان & معادل \cmd{a+x} است و به مالک، گروه و سایرین حق اجرا می‌دهد. \\
\addlinespace
\cmd{o-rw} & سلب حق خواندن و نوشتن از سایرین & دسترسی سایر کاربران (غیر از مالک و گروه) را به‌طور کامل برای خواندن و نوشتن مسدود می‌کند. \\
\addlinespace
\cmd{go=rw} & تنظیم مجوزهای خواندن و نوشتن برای گروه و سایرین & مجوز گروه و سایرین را دقیقاً به خواندن و نوشتن (\cmd{rw}) تنظیم می‌کند و در صورت وجود، حق اجرا (\cmd{x}) را از آن‌ها سلب می‌کند. \\
\addlinespace
\cmd{u+x,go=rx} & ترکیب چندین عملیات با کاما & با استفاده از ویرگول (\cmd{,}) دو عملیات مجزا در یک فرمان اجرا می‌شود: مالک حق اجرا می‌گیرد (\cmd{u+x}) و گروه و سایرین مجوزهای خواندن و اجرا (\cmd{rx}) را دریافت می‌کنند، اما حق نوشتن از آن‌ها سلب می‌شود (\cmd{go=rx}). \\
\bottomrule
\end{tabularx}
\end{table}

برای درک بهتر کاربرد روش نمادین، سناریوی زیر را در نظر بگیرید. فرض کنید فایلی به نام \file{report.txt} با مجوزهای پیش‌فرض \cmd{-rw-r--r--} (معادل ۶۴۴) دارید و قصد دارید بدون تغییر مجوزهای گروه و سایرین، تنها مجوز اجرا را به مالک فایل اضافه کنید. با استفاده از روش نمادین، این کار با یک دستور ساده انجام می‌شود:

\begin{latin}
\begin{lstlisting}
[me@linuxbox ~]$ ls -l report.txt
-rw-r--r-- 1 me me 0 2025-03-06 14:52 report.txt
[me@linuxbox ~]$ chmod u+x report.txt
[me@linuxbox ~]$ ls -l report.txt
-rwxr--r-- 1 me me 0 2025-03-06 14:52 report.txt
\end{lstlisting}
\end{latin}

در این مثال، عبارت نمادین \cmd{u+x} به‌روشنی نشان می‌دهد که فقط مجوز اجرا (\cmd{x}) به مالک (\cmd{u}) افزوده شده است، در حالی که سایر مجوزها (گروه و سایرین) کاملاً دست‌نخورده باقی می‌مانند. این سطح از کنترل دقیق، یکی از نقاط قوت روش نمادین است و آن را برای تنظیمات ظریف و محیط‌های اشتراکی بسیار مناسب می‌سازد.

برخی از مدیران سیستم استفاده از روش عددی (\en{Octal}) را به دلیل سرعت و فشردگی ترجیح می‌دهند، در حالی که برخی دیگر بازنمایی نمادین را به جهت خوانایی بالاتر، کاربردی‌تر می‌دانند. مزیت استراتژیک روش نمادین در این است که امکان اصلاح یک مجوز خاص را بدون دستکاری یا بازنویسی سایر دسترسی‌های موجود فراهم می‌کند؛ امری که انعطاف‌پذیری عملیاتی را به‌ویژه در محیط‌های اشتراکی به شکل قابل توجهی افزایش می‌دهد.

برای مطالعه جزئیات فنی دقیق‌تر و مشاهده فهرست کامل گزینه‌های عملیاتی، می‌توانید به مستندات رسمی سیستم از طریق فرمان \cmd{man chmod} مراجعه کنید.

هنگام کار با فرمان \cmd{chmod}، گاهی نیاز به اعمال تغییرات بر روی یک دایرکتوری و تمام محتویات درون آن به‌صورت یکجا داریم. این کار با استفاده از گزینهٔ \cmd{-R} (یا \cmd{--recursive}) امکان‌پذیر است. با این حال، به‌کارگیری این گزینه نیازمند دقت بالایی است؛ زیرا مجوزهای مناسب برای یک دایرکتوری (مانند مجوز \cmd{x} برای ورود به آن) ممکن است برای یک فایل معمولی (که نیازی به مجوز اجرا ندارد) کاملاً نامناسب و حتی مخاطره‌آمیز باشد.

\textbf{توصیهٔ امنیتی:} در بیشتر موارد، بهتر است به‌جای استفادهٔ سراسری از \cmd{chmod -R}، مجوزهای فایل‌ها و دایرکتوری‌ها را به‌صورت جداگانه و هدفمند تنظیم کنید. اگر ناگزیر به استفادهٔ بازگشتی هستید، پیش از اجرا، دامنهٔ تغییرات را به‌دقت بررسی کنید و در صورت امکان، از ترکیب \cmd{-R} با گزینه‌های نمادین (مانند \cmd{u+rwX} که در آن \cmd{X} تنها به دایرکتوری‌ها مجوز اجرا می‌دهد و نه فایل‌ها) برای کاهش ریسک استفاده کنید.

\subsection{پیکربندی سطوح دسترسی در محیط‌های گرافیکی (GUI)}

پس از درک عمیق مفاهیم فنی و منطق حاکم بر مجوزها، اکنون می‌توانیم ساختار پنجره‌های تنظیمات سطوح دسترسی را در محیط‌های دسکتاپ به شکلی دقیق‌تر تحلیل کنیم. در اکثر محیط‌های گرافیکی لینوکس نظیر \en{Files} (در محیط \en{GNOME}) یا \en{Dolphin} (در محیط \en{KDE})، با راست‌کلیک بر روی یک فایل یا دایرکتوری و انتخاب گزینه‌ی \en{Properties}، زبانه (\en{Tab}) ویژه‌ای برای مدیریت بصری سطوح دسترسی در اختیار کاربر قرار می‌گیرد.

این رابط‌های کاربری در واقع بازنمایی گرافیکی همان مفاهیم خواندن (\en{Read})، نوشتن (\en{Write}) و اجرا (\en{Execute}) هستند که برای سه رده‌ی مالک (\en{Owner})، گروه (\en{Group}) و سایرین (\en{Others}) تفکیک شده‌اند. این پنجره‌ها به عنوان جایگزینی ساده‌تر برای دستور \cmd{chmod} عمل می‌کنند و به کاربران اجازه می‌دهند بدون درگیری با سیستم هشت‌تایی یا فرمت‌های نمادین، سیاست‌های دسترسی را به سرعت اعمال کنند.

در محیطی مانند \en{GNOME}، تنظیمات به گونه‌ای طراحی شده‌اند که دسته‌های اصلی و مجوزهای متناظر با آن‌ها با وضوح کامل نمایش داده شوند. این رویکرد به‌ویژه برای کاربرانی که به تازگی به دنیای لینوکس وارد شده‌اند، راهکاری امن و بصری فراهم می‌کند تا از بروز خطاهای احتمالی در پیکربندی‌های متنی جلوگیری شود.

در ادامه، نمونه‌ای از پنجره‌ی تنظیمات مجوز در \en{GNOME} نمایش داده شده است که سه دسته‌ی اصلی و مجوزهای مربوط به هر یک را به وضوح نشان می‌دهد:

\bookfigure{img-02.png}{رابط گرافیکی مدیریت سطوح دسترسی در محیط GNOME}

در این پنجره عملیاتی، برای هر رده از کاربران (\en{Owner, Group, Others}) گزینه‌های مجزایی جهت پیکربندی مجوزهای خواندن (\en{Read})، نوشتن (\en{Write}) و اجرا (\en{Execute}) تعبیه شده است. این رابط کاربری به مدیران سیستم و کاربران اجازه می‌دهد بدون نیاز به استفاده از محیط ترمینال یا حفظ کردن کدهای هشت‌تایی، به سهولت وضعیت دسترسی‌ها را تغییر داده و امنیت داده‌های خود را مدیریت کنند.

این ساختار گرافیکی در واقع یک پوشش بصری بر روی همان منطق عددی لینوکس است؛ برای مثال، انتخاب حالت \en{``Read and Write''} در بخش \en{Owner}، معادل اجرای دستور \cmd{chmod 600} (در صورتی که سایرین دسترسی نداشته باشند) در لایه‌های زیرین سیستم‌عامل است.

\section{دستور umask: تنظیم مجوزهای پیش‌فرض}

دستور \cmd{umask} (مخفف \en{User File-Creation Mode Mask}) ابزاری است که تعیین می‌کند هنگام ایجاد فایل‌ها و دایرکتوری‌های جدید، کدام مجوزهای دسترسی به‌صورت پیش‌فرض از آن‌ها سلب شوند. این دستور از بازنمایی هشت‌تایی برای تعریف یک «ماسکِ بیتی» استفاده می‌کند؛ به عبارت دقیق‌تر، \cmd{umask} مشخص‌کننده‌ی مجوزهایی است که نباید روی فایل‌های جدید اعمال شوند.

برای درک بهتر سازوکار عملیاتی \cmd{umask}، سناریوی زیر را در نظر بگیرید:

\begin{latin}
\begin{lstlisting}
[me@linuxbox ~]$ rm -f foo.txt
[me@linuxbox ~]$ umask
0002
[me@linuxbox ~]$ > foo.txt
[me@linuxbox ~]$ ls -l foo.txt
-rw-rw-r-- 1 me me 0 2025-03-06 14:53 foo.txt
\end{lstlisting}
\end{latin}

در این مثال، ابتدا نسخه‌های پیشین فایل \file{foo.txt} حذف می‌شوند. سپس با اجرای دستور \cmd{umask} بدون هیچ پارامتری، مقدار فعلی ماسک سیستم را مشاهده می‌کنیم که در اینجا \cmd{0002} است (مقدار \cmd{0022} نیز در بسیاری از توزیع‌های لینوکس متداول است). پس از ایجاد یک فایل جدید، با استفاده از فرمان \cmd{ls -l} وضعیت مجوزهای آن را واکاوی می‌کنیم.

در منطق لینوکس، فایل‌ها به‌طور پیش‌فرض با سطح دسترسی 666 (خواندنی و نوشتنی برای همه) ایجاد می‌شوند، اما مقدار \cmd{umask} (در اینجا 002) از آن کسر می‌گردد. نتیجه این تفریق، مجوز 664 است که دقیقاً با خروجی \cmd{-rw-rw-r--} در مثال بالا مطابقت دارد.

همان‌طور که ملاحظه شد، مجوزهای پیش‌فرض فایل برابر با \cmd{-rw-rw-r--} تنظیم شده‌اند. برای تحلیل رابطه میان این دسترسی‌ها و مقدار \cmd{umask}، ابتدا باید از مجوزهای مبنا (\en{Base Permissions}) که سیستم برای فایل‌ها و دایرکتوری‌ها در نظر می‌گیرد آگاه باشیم؛ چرا که \cmd{umask} در واقع نقش یک «فیلتر» یا «کاهنده» را ایفا می‌کند که از این مقادیر پایه کسر می‌شود.

در مثال پیشین، دسترسی نهایی فایل جدید به صورت \cmd{-rw-rw-r--} بود؛ این یعنی مالک و گروه مجاز به خواندن و نوشتن هستند، اما سایر کاربران (\en{Others}) صرفاً حق خواندن دارند. علت سلب حق نوشتن از رده «سایرین»، دقیقاً به مقدار ماسک تنظیم شده در سیستم بازمی‌گردد.

برای کالبدشکافی این رفتار، سناریو را با یک ماسک خنثی تکرار می‌کنیم:

\begin{latin}
\begin{lstlisting}
[me@linuxbox ~]$ rm foo.txt
[me@linuxbox ~]$ umask 0000
[me@linuxbox ~]$ > foo.txt
[me@linuxbox ~]$ ls -l foo.txt
-rw-rw-rw- 1 me me 0 2025-03-06 14:58 foo.txt
\end{lstlisting}
\end{latin}

با تنظیم \cmd{umask} روی مقدار \cmd{0000} (وضعیت غیرفعال)، مشاهده می‌کنید که فایل جدید با دسترسی کاملِ خواندن و نوشتن (\cmd{-rw-rw-rw-}) برای تمامی رده‌ها ایجاد می‌شود. این نشان می‌دهد که سیستم به صورت پیش‌فرض تمایل دارد بالاترین سطح دسترسی غیرقابل اجرا (\en{Non-executable}) را اعمال کند، مگر آنکه \cmd{umask} مانع شود.

برای درک دقیقِ مکانیسم اثر \cmd{umask}، باید این مقدار را در سطح دودویی (\en{Binary}) بررسی کنیم. در واقع سیستم‌عامل یک عملیات منطقی \en{AND NOT} را بین مجوزهای پایه و ماسک انجام می‌دهد تا بیت‌های غیرمجاز را «خاموش» کند. در بخش بعد، با تحلیل مقدار \cmd{0002} خواهیم دید که چگونه بیتِ مربوط به «نوشتنِ دیگران» توسط این ماسک فیلتر می‌شود.

برای درک عمیق سازوکار \cmd{umask}، باید فرآیند را در سطح رایانش دودویی (\en{Binary}) واکاوی کرد. در واقع، \cmd{umask} یک «ماسکِ بیتی» (\en{Bitmask}) است که از طریق اعمال محدودیت بر مجوزهای پایه، سطح دسترسی نهایی را تعیین می‌کند.

\begin{latin}
\begin{lstlisting}
umask 0002
\end{lstlisting}
\end{latin}

به‌طور معمول، مجوز مبنا برای ایجاد یک فایل جدید در لینوکس، دسترسی کاملِ خواندن و نوشتن برای تمامی رده‌های کاربری (\en{Owner, Group, Others}) است که با نماد \cmd{-rw-rw-rw-} شناخته می‌شود. این وضعیت در سیستم دودویی به صورت مجموعه بیت‌های \cmd{110 110 110} بازنمایی می‌شود.

اکنون مقدار \cmd{umask 0002} را در نظر بگیرید. این عدد هشت‌تایی در مبنای دو به صورت زیر ترجمه می‌شود:

\begin{latin}
\begin{lstlisting}
000 000 010
\end{lstlisting}
\end{latin}

فرآیند تطبیق و حذف مجوزها به شرح زیر است:

\begin{quote}
حالت مبنای فایل: \lr{110\ 110\ 110}\\
ماسک اعمال‌شده: \lr{000\ 000\ 010}\\
خروجی نهایی: \lr{110\ 110\ 100}
\end{quote}

در این منطق محاسباتی، هرگاه بیتی در ماسک با مقدار «1» فعال باشد، مجوز متناظر با آن در فایل نهایی ابطال می‌گردد. در مثال فوق، تنها بیت فعال در ماسک (010) مربوط به مجوز نوشتن (\cmd{w}) برای رده سایرین (\en{Others}) است. در نتیجه، این حق دسترسی سلب شده و خروجی نهایی به صورت \cmd{-rw-rw-r--} تثبیت می‌شود.

\begin{latin}
\begin{lstlisting}
umask 0022
\end{lstlisting}
\end{latin}

در این مرحله، عملکرد مقدار رایج \cmd{umask 0022} را بررسی می‌کنیم. این مقدار در لایه منطقی و مبنای دودویی به صورت زیر تفسیر می‌شود:

\begin{latin}
\begin{lstlisting}
000 000 010 010
\end{lstlisting}
\end{latin}

فرآیند اعمال فیلتر بر روی مجوزهای پایه فایل به شرح زیر است:

\begin{quote}
حالت مبنای فایل: \lr{110\ 110\ 110}\\
ماسک اعمال‌شده: \lr{000\ 010\ 010}\\
خروجی نهایی: \lr{110\ 100\ 100}
\end{quote}

در این پیکربندی، ماسک دارای دو بیت فعال در موقعیت‌های استراتژیک است که حق نوشتن (\cmd{w}) را به‌طور هم‌زمان از گروه (\en{Group}) و سایرین (\en{Others}) سلب می‌کند. در نتیجه، فایل ایجاد شده تنها برای مالک قابل ویرایش بوده و سایر کاربران صرفاً مجاز به خواندن آن خواهند بود.

پس از انجام آزمایش‌های مختلف و درک چگونگی تغییر رفتار سیستم، می‌توانید با اجرای دستور زیر، تنظیمات نشست فعلی خود را به حالت استاندارد بازگردانید:

\begin{latin}
\begin{lstlisting}
[me@linuxbox ~]$ rm foo.txt; umask 0002
\end{lstlisting}
\end{latin}

در اکثر سناریوهای کاربری، نیازی به پیکربندی مجدد دستی \cmd{umask} نیست؛ چرا که مقادیر پیش‌فرض تعیین‌شده توسط توزیع‌های لینوکس (غالباً \cmd{0002} یا \cmd{0022}) توازن مناسبی میان امنیت و قابلیت استفاده برقرار کرده‌اند. با این حال، در زیرساخت‌های حساس و محیط‌های با سطح امنیت بالا، تنظیم دقیق \cmd{umask} یکی از گام‌های حیاتی برای اجرای سیاست‌های سخت‌گیرانه (\en{Hardening}) و کنترل دقیق دسترسی‌ها محسوب می‌شود.

\section{مجوزهای ویژه (Special Permissions)}

اگرچه در اکثر سناریوهای مدیریتی، سطوح دسترسی فایل‌ها در قالب اعداد سه‌رقمی هشت‌تایی (\en{Octal}) بیان می‌شوند، اما از منظر فنی، توصیف دقیق آن‌ها مستلزم استفاده از یک ساختار چهاررقمی است. رقم نخست در این سلسله‌مراتب، به سه «بیت ویژه» اختصاص دارد که اگرچه کمتر مورد استفاده قرار می‌گیرند، اما نقشی حیاتی در کنترل امنیت و مدیریت دسترسی‌ها در سطوح پیشرفته سیستم‌عامل ایفا می‌کنند.

\subsection*{بیت setuid (مقدار هشت‌تایی: 4000)}

هنگامی که بیت \cmd{setuid} (مخفف \en{Set User ID}) بر روی یک فایل اجرایی فعال می‌شود، «شناسه کاربری مؤثر» (\en{Effective UID}) فرآیند در حال اجرا، به شناسه مالک فایل تغییر می‌یابد. به عنوان مثال، چنانچه فایلی متعلق به کاربر \cmd{root} بوده و این بیت در آن فعال باشد، هر کاربر عادی با اجرای آن برنامه، به‌طور موقت از اختیارات و مجوزهای سطح \cmd{root} برخوردار خواهد شد. این مکانیسم برای عملکرد صحیح ابزارهایی نظیر \cmd{passwd} ضرورت فنی دارد؛ چرا که به کاربران اجازه می‌دهد فایل‌های حساس سیستمی (مانند \file{/etc/shadow}) را جهت تغییر گذرواژه خود اصلاح کنند؛ عملیاتی که در حالت عادی برای کاربران عادی مسدود است. با این وجود، به دلیل پتانسیل بالای این قابلیت در ایجاد رخنه‌های امنیتی، استفاده از آن باید به حداقل ممکن محدود شده و تحت نظارت دقیق امنیتی صورت گیرد.

\subsection*{بیت setgid (مقدار هشت‌تایی: 2000)}

بیت \cmd{setgid} (مخفف \en{Set Group ID}) دو کاربرد عملیاتی متمایز را در سطوح مختلف سیستم فایل ارائه می‌دهد:

\begin{itemize}
  \item \textbf{در فایل‌های اجرایی:} مشابه سازوکار \cmd{setuid}، این بیت «شناسه گروه مؤثر» (\en{Effective GID}) فرآیند را به شناسه گروهِ مالک فایل تغییر می‌دهد.
  \item \textbf{در دایرکتوری‌ها (کاربرد رایج):} این ویژگی نقشی کلیدی در مدیریت پروژه‌های گروهی ایفا می‌کند. چنانچه بیت \cmd{setgid} بر روی یک دایرکتوری تنظیم شود، هر فایل یا دایرکتوری جدیدی که در آن ایجاد گردد، به‌جای پذیرش گروهِ پیش‌فرضِ کاربر، به‌صورت خودکار مالکیت گروهیِ دایرکتوری والد را به ارث می‌برد. این قابلیت در محیط‌های اشتراکی تضمین می‌کند که اعضای یک تیم همواره به محتوای تولیدشده توسط سایر همکاران در آن دایرکتوری دسترسی داشته باشند.
\end{itemize}

\subsection*{بیت sticky (مقدار هشت‌تایی: 1000)}

بیت \cmd{sticky} پیشینه‌ای تاریخی در سیستم‌های یونیکس اولیه دارد که جهت بهینه‌سازی حافظه و ماندگاری برنامه‌ها (\en{Keep-in-memory}) استفاده می‌شد؛ اما در هسته لینوکس مدرن، این عملکرد برای فایل‌های معمولی منسوخ شده و نادیده گرفته می‌شود.

\begin{itemize}
  \item \textbf{در دایرکتوری‌ها:} در لینوکس امروزی، این بیت صرفاً برای دایرکتوری‌ها کاربرد دارد. فعال‌سازی بیت \cmd{sticky} بر روی یک دایرکتوری، یک لایه محافظتی ایجاد می‌کند که به موجب آن، تنها «مالک فایل»، «مالک دایرکتوری» یا «ابرکاربر (\cmd{root})» مجاز به حذف یا تغییر نام فایل‌های درون آن مسیر خواهند بود. این مکانیسم امنیتی به‌طور گسترده در دایرکتوری‌های عمومی مانند \file{/tmp} استفاده می‌شود تا از حذف یا دستکاری غیرمجاز فایل‌های یک کاربر توسط سایرین جلوگیری نموده و حریم خصوصی داده‌ها را در محیط‌های چندکاربره تضمین کند.
\end{itemize}

در ادامه، الگوهای عملیاتی فرمان \cmd{chmod} با بهره‌گیری از بازنمایی نمادین (\en{Symbolic}) جهت پیکربندی مجوزهای ویژه ارائه شده است:

برای فعال‌سازی بیت \cmd{setuid} بر روی یک فایل اجرایی، از پارامتر \cmd{u+s} استفاده می‌شود. این پیکربندی موجب می‌گردد برنامه با سطح دسترسی مالکِ فایل اجرا شود، حتی اگر فراخواننده‌ی آن کاربر دیگری باشد:

\begin{latin}
\begin{lstlisting}
chmod u+s program
\end{lstlisting}
\end{latin}

جهت تنظیم بیت \cmd{setgid} بر روی یک دایرکتوری، عبارت \cmd{g+s} به کار می‌رود. با این اقدام، تمامی فایل‌های نوظهور در این مسیر، به‌طور خودکار مالکیت گروهیِ دایرکتوریِ والد را کسب می‌کنند:

\begin{latin}
\begin{lstlisting}
chmod g+s dir
\end{lstlisting}
\end{latin}

برای اعمال بیت \cmd{sticky} بر روی یک دایرکتوری اشتراکی، از نماد \cmd{+t} استفاده می‌شود. این ویژگی امنیتی مانع از آن می‌شود که کاربران بتوانند فایل‌های متعلق به یکدیگر را حذف یا تغییر نام دهند:

\begin{latin}
\begin{lstlisting}
chmod +t dir
\end{lstlisting}
\end{latin}

در هنگام بازبینی خروجی فرمان \cmd{ls -l}، وضعیت این مجوزهای ویژه با نشانه‌های کاراکتری خاص در ستون دسترسی‌ها قابل تشخیص است:

\begin{itemize}
  \item \textbf{فایل مجهز به \cmd{setuid}:} وجود کاراکتر \cmd{s} در جایگاه مجوز اجرای مالک (\cmd{u})؛ مثال:
\begin{latin}
\begin{lstlisting}
-rwsr-xr-x
\end{lstlisting}
\end{latin}

  \item \textbf{دایرکتوری مجهز به \cmd{setgid}:} وجود کاراکتر \cmd{s} در جایگاه مجوز اجرای گروه (\cmd{g})؛ مثال:
\begin{latin}
\begin{lstlisting}
drwxrwsr-x
\end{lstlisting}
\end{latin}

  \item \textbf{دایرکتوری مجهز به \cmd{sticky bit}:} وجود کاراکتر \cmd{t} در آخرین جایگاه (مجوز اجرای سایرین یا \cmd{o})؛ مثال:
\begin{latin}
\begin{lstlisting}
drwxrwxrwt
\end{lstlisting}
\end{latin}
\end{itemize}

\begin{note}
\textbf{نکته:} اگر کاراکترهای مذکور به صورت حروف بزرگ (\cmd{S} یا \cmd{T}) ظاهر شوند، به این معناست که بیت ویژه فعال است اما مجوز اجرای پایه (\cmd{x}) برای آن رده تنظیم نشده است؛ در حالی که حروف کوچک نشان‌دهنده فعال بودن هم‌زمان بیت ویژه و مجوز اجرا هستند.
\end{note}

\section{تغییر هویت‌ها (Changing Identities)}

در اکوسیستم لینوکس، موقعیت‌های متعددی پیش می‌آید که در آن‌ها ناگزیر به احراز هویت تحت عنوانی دیگر هستیم. این ضرورت غالباً در هنگام اجرای فرآیندهای مدیریتی با استفاده از امتیازات ابرکاربر (\en{Superuser}) بروز می‌کند، هرچند ممکن است برای مقاصد عیب‌یابی و آزمایش، نیاز به شبیه‌سازی رفتار یک حساب کاربری عادی دیگر نیز وجود داشته باشد.

برای جابه‌جایی میان حساب‌های کاربری، سه راهکار عملیاتی وجود دارد:

\begin{itemize}
  \item \textbf{خروج و ورود مجدد (\en{Logout/Login}):} این روش، ضمن برخورداری از سادگی، به دلیل ماهیت زمان‌بر و قطع نشست فعلی، فاقد کارایی عملیاتی در محیط‌های حرفه‌ای است.
  \item \textbf{استفاده از فرمان \cmd{su}:} این ابزار امکان تغییر هویت به یک کاربر دیگر یا ایجاد یک پوسته (\en{Shell}) جدید با امتیازات آن کاربر را در همان نشست جاری فراهم می‌آورد. همچنین با استفاده از این فرمان می‌توان یک دستور واحد را با هویت کاربری متفاوت اجرا کرد.
  \item \textbf{استفاده از فرمان \cmd{sudo}:} این رویکرد به مدیران سیستم اجازه می‌دهد تا از طریق فایل پیکربندی \file{/etc/sudoers} به کاربران منتخب، مجوز اجرای دستورات خاص را با امتیازات سطح بالا اعطا کنند.
\end{itemize}

انتخاب میان \cmd{su} و \cmd{sudo} عمدتاً تابعی از سیاست‌های امنیتی و فلسفه طراحی توزیع لینوکس مورد استفاده است؛ با وجود در دسترس بودن هر دو ابزار، هر توزیع ممکن است یکی از این دو رویکرد را به عنوان استاندارد پیش‌فرض خود برگزیند.

در ادامه، جزئیات عملکرد دستور \cmd{su} را واکاوی خواهیم کرد، با این ملاحظه که در توزیع‌های مدرن لینوکس، تمایل به استفاده از این دستور در مقایسه با \cmd{sudo} کاهش یافته است.

\subsection{دستور su: اجرای پوسته با هویت جایگزین}

فرمان \cmd{su} (مخفف \en{Substitute User} یا \en{Switch User}) به‌منظور اجرای یک پوسته (\en{Shell}) با هویت کاربری دیگر طراحی شده است. ساختار نحوی این دستور به شرح زیر است:

\begin{latin}
\begin{lstlisting}
su [-[l]] [user]
\end{lstlisting}
\end{latin}

چنانچه از گزینه \cmd{-l} (یا به اختصار \cmd{-}) استفاده شود، پوسته جدید به عنوان یک «پوسته ورود» (\en{Login Shell}) برای کاربر هدف فراخوانی می‌گردد. این امر بدین معناست که محیط عملیاتی کاربر، شامل متغیرهای محیطی (\en{Environment Variables}) و مسیرهای اجرایی (\en{PATH})، به‌طور کامل بارگذاری شده و دایرکتوری کاری فعلی نیز به دایرکتوری خانگی (\en{Home Directory}) آن کاربر تغییر می‌یابد. در صورتی که نام کاربری در دستور قید نشود، \cmd{su} به‌طور پیش‌فرض هویت ابرکاربر (\cmd{root}) را هدف قرار می‌دهد.

به‌عنوان نمونه، برای دستیابی به پوسته \cmd{root} در سیستم‌هایی که حساب کاربری \cmd{root} دارای رمز عبور است، از دستور زیر استفاده می‌کنیم:

\begin{latin}
\begin{lstlisting}
[me@linuxbox ~]$ su -
Password:
[root@linuxbox ~]#
\end{lstlisting}
\end{latin}

پس از اجرای این فرمان، سیستم از شما رمز عبور کاربر \cmd{root} را مطالبه می‌کند. در صورت تایید، اعلان خط فرمان (\en{Prompt}) تغییر کرده و با نماد \cmd{\#} (به‌جای \cmd{\$}) نمایش داده می‌شود که نشان‌دهنده دسترسی با امتیازات ابرکاربر است. همچنین، مسیر جاری به \file{/root/} تغییر خواهد یافت.

برای خاتمه دادن به این نشست و بازگشت به هویت کاربری اولیه، کافی است دستور \cmd{exit} را اجرا نمایید:

\begin{latin}
\begin{lstlisting}
[root@linuxbox ~]# exit
[me@linuxbox ~]$
\end{lstlisting}
\end{latin}

\begin{note}
\textbf{نکته:} در بسیاری از توزیع‌های مدرن لینوکس (مانند اوبونتو)، حساب کاربری \cmd{root} به‌صورت پیش‌فرض فاقد رمز عبور است و ورود مستقیم از طریق \cmd{su -} میسر نیست. در این قبیل سیستم‌ها، معماران توزیع، بهره‌گیری از فرمان \cmd{sudo} را برای اجرای وظایف مدیریتی به عنوان رویکرد استاندارد و امن‌تر برگزیده‌اند.
\end{note}

در بسیاری از موارد، نیازی به باز کردن یک نشست تعاملی کامل نیست و صرفاً اجرای یک دستور واحد با امتیازات کاربری دیگر کفایت می‌کند. برای این منظور، فرمان \cmd{su} گزینه \cmd{-c} (مخفف \en{command}) را ارائه می‌دهد. ساختار نحوی این دستور به شرح زیر است:

\begin{latin}
\begin{lstlisting}
su -c 'command'
\end{lstlisting}
\end{latin}

در این ساختار، دستور مورد نظر باید درون علامت کوتیشن (\en{Quotation}) محصور شود. این نکتهٔ نحوی تضمین می‌کند که دستور توسط پوسته (\en{Shell}) جدید و تحت هویت کاربر جایگزین پردازش شود، نه در محیط پوسته فعلی شما.

\textbf{مثال عملیاتی:} به‌عنوان نمونه، برای فهرست کردن محتویات دایرکتوری \file{/root/} که به‌طور معمول برای کاربران عادی غیرقابل‌دسترس است، می‌توان از این روش بهره جست:

\begin{latin}
\begin{lstlisting}
[me@linuxbox ~]$ su -c 'ls -l /root/*'
Password:
-rw------- 1 root root 754 2007-08-11 03:19 /root/anaconda-ks.cfg

/root/Mail:
total 0
[me@linuxbox ~]$
\end{lstlisting}
\end{latin}

در این سناریو، فرمان \cmd{ls} با پارامتر \cmd{-l} بر روی محتویات مسیر \file{/root/} و با اختیارات ابرکاربر (\en{Superuser}) اجرا شده است. پس از تایید رمز عبور، خروجی دستور در خروجی استاندارد چاپ شده و بلافاصله کنترل نشست به پوسته اولیه بازمی‌گردد.

این روش از منظر امنیتی بر باز کردن یک نشست کامل ارجحیت دارد، زیرا ریسکِ باز ماندن یک پوسته با امتیازات سطح بالا (\en{Root Shell}) را به حداقل می‌رساند.

\subsection{دستور sudo: اجرای مدیریت‌شده دستورات با امتیازات ارتقایافته}

فرمان \cmd{sudo} (مخفف \en{Superuser Do} یا \en{Substitute User Do}) ابزاری است که به کاربران اجازه می‌دهد دستورات را با امتیازات امنیتی کاربر دیگر (غالباً ابرکاربر یا \cmd{root}) اجرا کنند. این ابزار در مقایسه با \cmd{su} از انعطاف‌پذیری و قابلیت‌های کنترلی بسیار دقیق‌تری برخوردار است. مدیر سیستم می‌تواند با پیکربندی فایل \file{/etc/sudoers} مشخص کند که هر کاربر دقیقاً مجاز به اجرای چه مجموعه‌ای از دستورات با چه سطح دسترسی است.

یکی از مزایای بنیادین \cmd{sudo} از منظر امنیتی، عدم نیاز به اشتراک‌گذاری رمز عبور ابرکاربر است. کاربران در این روش تنها با وارد کردن رمز عبور شخصی خود، هویتشان را احراز می‌کنند. این رویکرد ریسک افشای رمز عبور اصلی سیستم (\en{Root Password}) را به‌طور جدی کاهش می‌دهد.

برای واکاوی عملکرد عملیاتی \cmd{sudo}، سناریوی اجرای یک اسکریپت پشتیبان‌گیری (\en{Backup}) را که مستلزم امتیازات مدیریتی است، در نظر بگیرید:

\begin{latin}
\begin{lstlisting}
[me@linuxbox ~]$ sudo backup_script
Password:
System Backup Starting...
\end{lstlisting}
\end{latin}

پس از درج دستور \cmd{sudo backup\_script} و فشردن کلید اینتر، سیستم رمز عبورِ کاربر فعلی را مطالبه می‌کند. پس از تایید هویت، دستور با امتیازات لازم فراخوانی شده و فرآیند پشتیبان‌گیری آغاز می‌گردد. این مکانیسم، امکان واگذاری اختیار به کاربران مورد اعتماد را بدون نیاز به واگذاری کنترل کامل سیستم فراهم می‌آورد.

یکی دیگر از تمایزهای بنیادین میان \cmd{sudo} و \cmd{su} در این است که \cmd{sudo} به‌طور پیش‌فرض پوسته (\en{Shell}) جدیدی ایجاد نکرده و محیط عملیاتی کاربر مقصد (مانند متغیرهای محیطی) را بارگذاری نمی‌کند. این ویژگی موجب می‌شود که برخلاف دستور \cmd{su -c} نیازی به محصور کردن فرامین در علامت کوتیشن (\en{Quotation}) نباشد. البته این رفتارِ پیش‌فرض با استفاده از گزینه‌های اختصاصی قابل تغییر است.

به‌عنوان نمونه، با بهره‌گیری از گزینه \cmd{-i} (مخفف \en{Interactive})، می‌توانید یک نشست تعاملی با امتیازات ابرکاربر (\en{Superuser}) ایجاد کنید که از نظر عملکردی و بارگذاری پروفایل‌های کاربری، مشابه دستور \cmd{su -} عمل می‌کند:

\begin{latin}
\begin{lstlisting}
[me@linuxbox ~]$ sudo -i
\end{lstlisting}
\end{latin}

در این حالت، متغیرهای محیطی نظیر \cmd{PATH} و \cmd{HOME} متناسب با حساب کاربری \cmd{root} تنظیم شده و دایرکتوری کاری شما نیز به دایرکتوری خانگی ابرکاربر تغییر می‌یابد. این روش، راهکاری امن و استاندارد برای مدیریت طولانی‌مدت سیستم بدون نیاز به فعال‌سازی رمز عبور مستقیم برای حساب کاربری \cmd{root} است.

به منظور استخراج و مشاهده فهرست دقیقِ فرامین و اختیاراتی که از طریق \cmd{sudo} به حساب کاربری شما تفویض شده است، می‌توانید از گزینه \cmd{-l} (مخفف \en{List}) استفاده نمایید:

\begin{latin}
\begin{lstlisting}
[me@linuxbox ~]$ sudo -l
User me may run the following commands on this host:
    (ALL) ALL
\end{lstlisting}
\end{latin}

خروجی فوق بیانگر آن است که کاربر \en{me} مجاز است تمامی دستورات را با امتیازات هر کاربری (\cmd{ALL}) بر روی این ماشین اجرا کند. این سطح از دسترسی، پیکربندی پیش‌فرض در بسیاری از توزیع‌های مدرن لینوکس محسوب می‌شود و به کاربرانِ معتمد اجازه می‌دهد تا فرآیندهای مدیریتی و نگهداری سیستم را بدون محدودیت‌های ساختاری و با سهولت عملیاتی به انجام رسانند.

\subsection{مدیریت امتیازات با sudo در توزیع‌های مدرن لینوکس}

در سیستم‌عامل‌های چندکاربره نظیر لینوکس، مدیریت عملیاتی که مستلزم امتیازات ابرکاربر (\en{Superuser}) است — از جمله نصب بسته‌های نرم‌افزاری، به‌روزرسانی هسته سیستم یا اصلاح فایل‌های پیکربندی حساس — یک چالش امنیتی بنیادین محسوب می‌شود. در گذشته، بسیاری از توزیع‌های لینوکس بر اساس سنت‌های کلاسیک یونیکس، از فرمان \cmd{su} برای این مقاصد بهره می‌بردند. این رویکرد کاربران را به تعامل مستقیم و مستمر با حساب کاربری \cmd{root} ترغیب می‌کرد که به مرور زمان، سطح آسیب‌پذیری سیستم را افزایش می‌داد.

با ظهور توزیع اوبونتو، الگوی جدیدی معرفی شد که امروزه به استاندارد غالب در توزیع‌های مدرن مبدل گشته است. در این مدل نوین، به‌صورت پیش‌فرض:

\begin{itemize}
  \item \textbf{غیرفعال‌سازی رمز عبور حساب \cmd{root}:} با حذف رمز عبور برای ابرکاربر، امکان ورود مستقیم (\en{Login}) به این حساب به‌طور کامل مسدود شده است.
  \item \textbf{محوریت فرمان \cmd{sudo} در واگذاری اختیارات:} کاربر اولیه سیستم به‌طور خودکار در گروه کاربری \cmd{sudo} (یا در برخی توزیع‌ها گروه \cmd{wheel}) قرار می‌گیرد و بدین ترتیب مجاز است فرامین مدیریتی را از طریق پروتکل \cmd{sudo} به اجرا درآورد.
\end{itemize}

این متدولوژی، ضریب امنیت سیستم را به‌طور قابل‌توجهی ارتقا می‌دهد؛ چرا که کاربران تنها در موارد اضطراری و برای اجرای فرامین مشخص، امتیازات مدیریتی را احراز می‌کنند. این تمایز در فلسفه امنیتی، لینوکس را در مقایسه با سیستم‌عامل‌هایی که به‌صورت پیش‌فرض دسترسی‌های مدیریتی گسترده‌ای به کاربران عادی اعطا می‌کنند، در برابر بدافزارها و کدهای مخرب بسیار مقاوم‌تر ساخته است.

\subsection{دستور chown: تغییر مالک و گروه فایل}

فرمان \cmd{chown} (برگرفته از عبارت \en{change owner}) یکی از ابزارهای بنیادین در سیستم‌عامل لینوکس است که وظیفه تغییر «مالکیت کاربری» (\en{Owner}) و «مالکیت گروهی» (\en{Group}) فایل‌ها یا دایرکتوری‌ها را بر عهده دارد. با توجه به اینکه اصلاح ساختار مالکیت، عملیاتی در سطح هسته و امنیت سیستم محسوب می‌شود، اجرای این دستور مستلزم برخورداری از امتیازات ابرکاربر (\en{Superuser}) است.

ساختار نحوی (\en{Syntax}) این فرمان به شرح زیر است:

\begin{latin}
\begin{lstlisting}
chown [owner][:group] file...
\end{lstlisting}
\end{latin}

آرگومان نخست در این ساختار، تعیین‌کننده هویت مالک و گروه جدید است. در این چارچوب، سیستم به شما اجازه می‌دهد تا مالک، گروه، یا هر دو را به‌صورت هم‌زمان بازتعریف کنید.

\begin{table}[htbp]
\centering
\caption{آرگومان‌های کاربردی دستور \cmd{chown}}
\label{tab:chown-args}
\begin{tabularx}{\textwidth}{l X}
\toprule
\textbf{آرگومان ورودی} & \textbf{پیامد عملیاتی} \\
\midrule
\cmd{bob} & مالکیت فایل از کاربر فعلی به کاربر \en{bob} منتقل می‌شود. \\
\addlinespace
\cmd{bob:users} & مالکیت فایل به کاربر \en{bob} واگذار شده و گروهِ مالک نیز به \en{users} تغییر می‌یابد. \\
\addlinespace
\cmd{:admins} & تنها مالکیت گروهی به \en{admins} تغییر می‌کند و مالک اصلی فایل بدون تغییر باقی می‌ماند. \\
\addlinespace
\cmd{bob:} & مالکیت به \en{bob} منتقل شده و گروهِ مالک نیز به‌طور خودکار به «گروه اصلی» (\en{Login Group}) کاربر \en{bob} تغییر می‌یابد. \\
\bottomrule
\end{tabularx}
\end{table}

در سناریوی پیش‌رو، کاربر \en{janet} (دارای امتیازات مدیریتی) قصد دارد فایلی تحت عنوان \file{myfile.txt} را از دایرکتوری خانگی خود به مسیر خانگی کاربر \en{tony} (کاربر عادی) منتقل نماید. هدف این است که پس از انتقال، \en{tony} دسترسی کامل برای ویرایش فایل را داشته باشد. برای این منظور، فرآیند زیر طی می‌شود:

\begin{latin}
\begin{lstlisting}
[janet@linuxbox ~]$ sudo cp myfile.txt ~tony
Password:
[janet@linuxbox ~]$ sudo ls -l ~tony/myfile.txt
-rw-r--r-- 1 root root root 2025-03-20 14:30 /home/tony/myfile.txt
[janet@linuxbox ~]$ sudo chown tony: ~tony/myfile.txt
[janet@linuxbox ~]$ sudo ls -l ~tony/myfile.txt
-rw-r--r-- 1 tony tony tony 2025-03-20 14:30 /home/tony/myfile.txt
\end{lstlisting}
\end{latin}

تحلیل فنی مراحل فوق به شرح زیر است:

\begin{itemize}
  \item \textbf{تغییر مالکیت ناخواسته هنگام کپی:} ابتدا \en{janet} با بهره‌گیری از فرمان \cmd{cp} و امتیازات \cmd{sudo} فایل را کپی می‌کند. از آنجا که عملیات کپی تحت هویت ابرکاربر انجام شده است، مالکیت فایل جدید به‌صورت خودکار به کاربر و گروه \cmd{root} تعلق می‌گیرد.
  \item \textbf{اصلاح سلسله‌مراتب مالکیت:} در مرحله بعد، \en{janet} با اجرای فرمان \cmd{sudo chown tony:} مالکیت فایل را به \en{tony} واگذار می‌کند. درج علامت دونقطه (\cmd{:}) پس از نام کاربری، یک تکنیک کاربردی است که موجب می‌شود «گروه مالک» نیز به‌طور خودکار به گروه اصلی (\en{Primary Group}) کاربر \en{tony} تغییر یابد.
  \item \textbf{تأیید نهایی:} خروجی دومِ فرمان \cmd{ls -l} به‌وضوح نشان می‌دهد که هر دو فیلد مالک و گروه به \en{tony} تخصیص یافته‌اند، که این امر امکان مدیریت کامل فایل را برای وی فراهم می‌سازد.
\end{itemize}

شایان ذکر است که پس از نخستین احراز هویت در فرمان اول، سیستم برای اجرای دستور دوم مطالبه رمز عبور نمی‌کند. این پدیده به دلیل وجود سیاست تداوم نشست (\en{Ticket Timestamp}) در پیکربندی \cmd{sudo} است که به کاربر اجازه می‌دهد برای بازه‌ای کوتاه (معمولاً ۵ تا ۱۵ دقیقه)، دستورات بعدی را بدون تکرار فرآیند احراز هویت اجرا کند.

\subsection{دستور chgrp: تغییر گروه مالکیت}

در نسخه‌های آغازین سیستم‌عامل یونیکس، ابزار \cmd{chown} منحصراً برای تغییر کاربرِ مالک (\en{Owner}) طراحی شده بود و قابلیتی برای مدیریت گروه‌ها نداشت. به همین منظور، فرمان مجزای \cmd{chgrp} (برگرفته از عبارت \en{change group}) برای مدیریت اختصاصی گروه‌های مالکیت توسعه یافت.

عملکرد \cmd{chgrp} شباهت ساختاری بسیاری به \cmd{chown} دارد، با این تفاوت که دامنه عملیاتی آن محدود به تغییر «گروه مالکیت» فایل‌ها یا دایرکتوری‌هاست. با وجود آنکه در توزیع‌های مدرن لینوکس، فرمان \cmd{chown} تکامل یافته و امکان تغییر هم‌زمان کاربر و گروه را (با استفاده از نحو \cmd{user:group}) فراهم کرده است، اما \cmd{chgrp} همچنان به عنوان یک ابزار تخصصی و مستقل در جعبه‌ابزار مدیریت سیستم (\en{Sysadmin}) جایگاه خود را حفظ کرده است.

\section{تمرین: ایجاد دایرکتوری اشتراکی}

پس از درک مفاهیم تئوریک مجوزها، اکنون این آموخته‌ها را در قالب یک سناریوی کاربردی پیاده‌سازی می‌کنیم: «ایجاد یک دایرکتوری اشتراکی». در این مثال، از دو کاربر با نام‌های \en{janet} (دارای امتیازات ابرکاربر یا \en{Superuser}) و \en{tony} استفاده خواهیم کرد. هدف نهایی، پیکربندی یک دایرکتوری مشترک جهت تبادل فایل‌های صوتی است، به‌گونه‌ای که هر دو کاربر دسترسی‌های لازم را دارا باشند.

این فرآیند عملیاتی طی دو مرحله کلیدی اجرا می‌شود:

\textbf{ایجاد گروه جدید:} در گام نخست، باید یک گروه تحت عنوان \cmd{music} تعریف کنیم تا به عنوان محلی برای عضویت هر دو کاربر عمل کند. این عمل از طریق فرمان \cmd{groupadd} صورت می‌پذیرد:

\begin{latin}
\begin{lstlisting}
[janet@linuxbox ~]$ sudo groupadd music
\end{lstlisting}
\end{latin}

\textbf{افزودن کاربران به گروه:} در گام دوم، کاربران \en{janet} و \en{tony} را به عضویت گروه \cmd{music} در می‌آوریم. برای این منظور از فرمان \cmd{usermod} بهره می‌گیریم:

\begin{latin}
\begin{lstlisting}
[janet@linuxbox ~]$ sudo usermod -a -G music janet
[janet@linuxbox ~]$ sudo usermod -a -G music tony
\end{lstlisting}
\end{latin}

تحلیل گزینه‌های به‌کار رفته در دستور \cmd{usermod}:

\begin{itemize}
  \item \cmd{-a} (مخفف \en{append}): این گزینه تضمین می‌کند که کاربر به گروه (یا گروه‌های) جدید افزوده شود، بدون آنکه عضویت او در گروه‌های فعلی لغو گردد.
  \item \cmd{-G} (مخفف \en{groups}): این پارامتر برای تعیین نام گروه‌های ثانویه‌ای که کاربر باید به آن‌ها ملحق شود، استفاده می‌گردد.
\end{itemize}

با اجرای موفقیت‌آمیز این دستورات، گروه \cmd{music} در لایه سیستم‌عامل تعریف شده و هر دو کاربر به عضویت آن در می‌آیند. این تغییرات به‌صورت دائمی در فایل پیکربندی \file{/etc/group} ثبت و ذخیره می‌شود.

در مرحله بعد، \en{janet} دایرکتوری مورد نظر را جهت استقرار فایل‌های موسیقی ایجاد می‌کند. از آنجایی که مسیر انتخابی خارج از دایرکتوری خانگی (\en{Home Directory}) اوست، اجرای این عملیات مستلزم برخورداری از امتیازات ابرکاربر است؛ لذا \en{janet} از دستور \cmd{sudo} بهره می‌گیرد:

\begin{latin}
\begin{lstlisting}
[janet@linuxbox ~]$ sudo mkdir /usr/local/share/Music
Password:
\end{lstlisting}
\end{latin}

پس از اتمام فرآیند ایجاد، وضعیت مالکیت و سطوح دسترسی پیش‌فرض دایرکتوری به شرح زیر است:

\begin{latin}
\begin{lstlisting}
[janet@linuxbox ~]$ ls -ld /usr/local/share/Music
drwxr-xr-x 2 root root 4096 2025-03-21 18:05 /usr/local/share/Music
\end{lstlisting}
\end{latin}

خروجی فوق نشان می‌دهد که دایرکتوری مذکور تحت مالکیت کاربر \cmd{root} و گروه \cmd{root} قرار دارد و با مجوزهای استاندارد 755 (یعنی \cmd{drwxr-xr-x}) ایجاد شده است. این پیکربندی به اعضای گروه و سایر کاربران (\en{Others}) تنها اجازه «خواندن» و «اجرا» را می‌دهد. برای تبدیل این مسیر به یک فضای اشتراکی (\en{Shared Directory})، \en{janet} باید مالکیت گروهی و مجوزهای دسترسی را به‌گونه‌ای بازتنظیم کند که اعضای گروه \cmd{music} قابلیت «نوشتن» و ایجاد فایل را نیز داشته باشند:

\begin{latin}
\begin{lstlisting}
[janet@linuxbox ~]$ sudo chown :music /usr/local/share/Music
[janet@linuxbox ~]$ sudo chmod 2775 /usr/local/share/Music
\end{lstlisting}
\end{latin}

پس از اعمال تغییرات، وضعیت نهایی دایرکتوری به صورت زیر خواهد بود:

\begin{latin}
\begin{lstlisting}
[janet@linuxbox ~]$ ls -ld /usr/local/share/Music
drwxrwsr-x 2 root music 4096 2025-03-21 18:05 /usr/local/share/Music
\end{lstlisting}
\end{latin}

در این مرحله، استفاده از عدد 2 در ابتدای مجوز 2775 باعث فعال‌سازی بیت \cmd{setgid} (خروجی \cmd{s} در بخش مجوز گروه) شده است. این ویژگی تضمین می‌کند که تمامی فایل‌های جدید ایجاد شده در این دایرکتوری، به‌طور خودکار مالکیت گروهی \cmd{music} را به ارث ببرند.

در این فرآیند عملیاتی:

\begin{itemize}
  \item \textbf{فرمان \cmd{chown}:} با اجرای دستور:
\begin{latin}
\begin{lstlisting}
[janet@linuxbox ~]$ sudo chown :music /usr/local/share/Music
\end{lstlisting}
\end{latin}
توسط \en{janet}، مالکیت گروهی دایرکتوری به گروه \cmd{music} واگذار می‌شود. این اقدام باعث می‌شود تا تمامی دسترسی‌های تعریف‌شده برای سطح گروه، مستقیماً به اعضای گروه \cmd{music} (یعنی \en{janet} و \en{tony}) اختصاص یابد.

  \item \textbf{فرمان \cmd{chmod}:} در گام بعد، با استفاده از دستور:
\begin{latin}
\begin{lstlisting}
[janet@linuxbox ~]$ sudo chmod 2775 /usr/local/share/Music
[janet@linuxbox ~]$ ls -ld /usr/local/share/Music
drwxrwsr-x 2 root music 4096 2025-03-21 18:05 /usr/local/share/Music
\end{lstlisting}
\end{latin}
سطوح دسترسی دایرکتوری بازتنظیم می‌شود. مقدار عددی 2775 در این سناریو نقشی حیاتی ایفا می‌کند:
\begin{itemize}
  \item \textbf{رقم 2 (\cmd{setgid}):} این رقم بیت \cmd{setgid} را فعال می‌کند. هنگامی که این ویژگی بر روی یک دایرکتوری تنظیم شود، تمامی فایل‌ها و دایرکتوری‌های فرعی جدیدی که درون آن ایجاد می‌شوند، به‌جای ارث‌بری از گروه پیش‌فرضِ کاربرِ ایجادکننده، به‌طور خودکار مالکیت گروهی دایرکتوری والد (در اینجا \cmd{music}) را به ارث می‌برند. این قابلیت، پیش‌نیاز اصلی برای یک اشتراک‌گذاری کارآمد و بی‌نقص است.
  \item \textbf{رقم 775:} این بخش از کد، مجوزهای «خواندن، نوشتن و اجرا» (\cmd{rwx}) را برای کاربر مالک (\cmd{root}) و گروه مالکیت (\cmd{music}) تعیین کرده و برای سایر کاربران سیستم (\en{Others})، تنها مجوزهای «خواندن و اجرا» (\cmd{r-x}) را لحاظ می‌کند.
\end{itemize}
\end{itemize}

حاصل این تنظیمات، دایرکتوری \file{/usr/local/share/Music/} است که تحت مالکیت کاربر \cmd{root} قرار دارد اما مدیریت گروهی آن در اختیار \cmd{music} است. به علت فعال بودن بیت \cmd{setgid} و اختصاص مجوزهای \cmd{rwx} به گروه، اعضای گروه \cmd{music} (\en{janet} و \en{tony}) قادر خواهند بود فایل‌های جدید را ایجاد و ویرایش کنند. در مقابل، سایر کاربران سیستم به دلیل محدودیت دسترسی در سطح \cmd{r-x} تنها مجاز به مشاهده و پیمایش محتویات این دایرکتوری خواهند بود.

با وجود تنظیمات قبلی، کماکان یک مانع فنی در مسیر همکاری مشترک وجود دارد: مقدار پیش‌فرض \cmd{umask} در این سیستم بر روی \cmd{0022} تنظیم شده است. این پیکربندی به این معناست که هر فایل جدیدی که توسط \en{janet} یا \en{tony} ایجاد شود، به‌طور خودکار فاقد مجوز «نوشتن» برای سایر اعضای گروه خواهد بود. به بیان ساده‌تر، چنانچه \en{janet} فایلی را ایجاد کند، \en{tony} علیرغم عضویت در گروه مشترک، اجازه ویرایش آن را نخواهد داشت. برای رفع این محدودیت، لازم است مقدار \cmd{umask} برای هر دو کاربر به \cmd{0002} تغییر یابد؛ این اقدام تضمین می‌کند که فایل‌های جدید، مجوزهای نگارش گروهی را حفظ کرده و فرآیند همکاری بدون اختلال تداوم یابد.

\en{janet} با بازتنظیم \cmd{umask} خود بر روی مقدار \cmd{0002} اطمینان حاصل می‌کند که هر فایل یا دایرکتوری جدیدی که در این فضای اشتراکی ایجاد می‌شود، دارای مجوزهایی است که به طور صریح اجازه «نوشتن» و «اصلاح» را به سایر اعضای گروه \cmd{music} اعطا می‌کند.

پس از بازتنظیم مقدار \cmd{umask}، کاربر \en{janet} اقدام به ایجاد یک فایل و یک دایرکتوری جدید در مسیر \file{/usr/local/share/Music/} می‌کند تا صحت پیکربندی‌ها را بررسی نماید:

\begin{latin}
\begin{lstlisting}
[janet@linuxbox ~]$ umask 0002
[janet@linuxbox ~]$ > /usr/local/share/Music/test_file
[janet@linuxbox ~]$ mkdir /usr/local/share/Music/test_dir
[janet@linuxbox ~]$ ls -l /usr/local/share/Music
drwxrwsr-x 2 janet music 4096 2025-03-24 20:24 test_dir
-rw-rw-r-- 1 janet music    0 2025-03-24 20:22 test_file
[janet@linuxbox ~]$
\end{lstlisting}
\end{latin}

با تحلیل خروجی \cmd{ls -l} مشاهده می‌شود که فایل \file{test\_file} با مجوز \cmd{-rw-rw-r--} و دایرکتوری \file{test\_dir} با مجوز \cmd{drwxrwsr-x} ایجاد شده‌اند. این سطوح دسترسی، حق «خواندن و نوشتن» کامل را برای مالک (\en{janet}) و اعضای گروه (\cmd{music}) تضمین می‌کنند.

در این وضعیت، به دلیل ارث‌بری صحیح مالکیت گروهی (به واسطه بیت \cmd{setgid}) و آزادسازی مجوز نوشتن گروهی (به واسطه \cmd{umask 0002})، کاربر \en{tony} نیز به عنوان عضو گروه \cmd{music} می‌تواند بدون محدودیت، فایل‌ها و دایرکتوری‌های ایجاد شده توسط \en{janet} را ویرایش یا مدیریت کند.

\begin{note}
\textbf{نکته:} تغییر \cmd{umask} از طریق خط فرمان، صرفاً به‌صورت موقت اعمال شده و تنها برای نشست جاری (\en{Session}) معتبر است. به محض خروج کاربر از سیستم یا بستن محیط ترمینال، این تنظیمات به مقادیر پیش‌فرض بازخواهند گشت.
\end{note}

جهت دائمی کردن این تغییرات و اطمینان از اعمال خودکار آن‌ها در هر بار ورود به سیستم، بایستی دستور مربوطه را در فایل‌های پیکربندی پوسته کاربر (مانند \file{.bashrc} یا \file{.profile}) درج نمود. جزئیات دقیق و نحوه مدیریت این فایل‌های پیکربندی در فصل ۱۱ به‌طور تفصیلی مورد بررسی قرار خواهد گرفت.

\subsection{دستور passwd: تغییر گذرواژه کاربر}

آخرین مبحثی که در این فصل بررسی می‌کنیم، مدیریت امنیت و تنظیمات گذرواژه (\en{Password}) کاربران است. برای پیکربندی یا تغییر رمز عبور شخصی و یا مدیریت رمز عبور سایر کاربران سیستم (در صورت برخورداری از امتیازات ابرکاربر)، از ابزار \cmd{passwd} استفاده می‌شود. نحو (\en{Syntax}) کلی این دستور به صورت زیر است:

\begin{latin}
\begin{lstlisting}
passwd [user]
\end{lstlisting}
\end{latin}

جهت تغییر گذرواژه حساب کاربری فعلی، کافی است دستور \cmd{passwd} را به تنهایی اجرا کنید. سیستم در ابتدا گذرواژه جاری و سپس گذرواژه جدید را از شما درخواست خواهد کرد:

\begin{latin}
\begin{lstlisting}
[me@linuxbox ~]$ passwd
(current) UNIX password:
New UNIX password:
\end{lstlisting}
\end{latin}

در صورتی که پارامتر \file{user} وارد نشود، دستور \cmd{passwd} به‌طور پیش‌فرض بر روی حساب کاربری فعلی اعمال می‌گردد. شایان ذکر است که تغییر گذرواژه سایر کاربران تنها با دسترسی ریشه (\en{Root}) یا استفاده از امتیازات \cmd{sudo} امکان‌پذیر است.

فرمان \cmd{passwd} به منظور ارتقای امنیت سیستم، معیارهای سخت‌گیرانه‌ای را برای احراز «گذرواژه‌های مستحکم» (\en{Strong Passwords}) اعمال می‌کند. این ابزار از پذیرش رمزهایی که دارای ویژگی‌های زیر باشند، خودداری خواهد کرد:

\begin{itemize}
  \item طول کاراکترهای آن‌ها بسیار کوتاه باشد.
  \item شباهت ساختاری زیادی با گذرواژه پیشین داشته باشند.
  \item در دسته‌بندی واژگان متداول لغت‌نامه (\en{Dictionary Words}) قرار بگیرند.
  \item بر اساس الگوهای ساده و قابل حدس طراحی شده باشند.
\end{itemize}

در مثال عملی زیر، فرآیند تلاش برای انتخاب یک رمز عبور ضعیف و واکنش‌های حفاظتی سیستم به تصویر کشیده شده است:

\begin{latin}
\begin{lstlisting}
[me@linuxbox ~]$ passwd
(current) UNIX password:
New UNIX password:
BAD PASSWORD: is too similar to the old one
New UNIX password:
BAD PASSWORD: it is WAY too short
New UNIX password:
BAD PASSWORD: it is based on a dictionary word
\end{lstlisting}
\end{latin}

همان‌طور که در خروجی بالا مشهود است، سیستم به صورت هوشمند از پذیرش گذرواژه‌های ناامن خودداری کرده و هشدارهای امنیتی متناسب با هر خطا را صادر می‌کند. این پیام‌ها به ترتیب بیانگر «شباهت بیش از حد به رمز پیشین»، «کوتاهی مفرط طول کاراکترها» و «استفاده از واژگان رایج لغت‌نامه» هستند که همگی از منظر امنیتی، آسیب‌پذیری سیستم را افزایش می‌دهند.

چنانچه به امتیازات ابرکاربر (\en{Superuser}) دسترسی داشته باشید، می‌توانید با درج نام کاربری به‌عنوان آرگومان در فرمان \cmd{passwd} اقدام به تغییر گذرواژه سایر کاربران نمایید.

\textbf{مثال عملی:}

\begin{latin}
\begin{lstlisting}
[me@linuxbox ~]$ sudo passwd tony
\end{lstlisting}
\end{latin}

این دستور به شما این امکان را می‌دهد تا گذرواژه کاربر \en{tony} را به یک رشته متنی جدید و مستحکم تغییر دهید. فراتر از تغییر ساده رمز عبور، فرمان \cmd{passwd} مجموعه‌ای از گزینه‌های عملیاتی پیشرفته را در اختیار مدیر سیستم قرار می‌دهد؛ قابلیت‌هایی نظیر:

\begin{itemize}
  \item قفل کردن حساب کاربری (\en{Locking}): غیرفعال‌سازی موقت دسترسی کاربر به سیستم.
  \item انقضای اجباری گذرواژه (\en{Expiring}): الزام کاربر به تغییر رمز عبور در اولین ورود بعدی.
  \item سایر تنظیمات امنیتی: مدیریت دوره‌های زمانی معتبر بودن گذرواژه‌ها.
\end{itemize}

برای بررسی جامع تمامی گزینه‌ها و قابلیت‌های این ابزار، مطالعه صفحات راهنما از طریق دستور \cmd{man passwd} توصیه می‌شود.

فرامین \cmd{passwd}، \cmd{groupadd} و \cmd{usermod} تنها بخشی از بسته نرم‌افزاری جامع \cmd{shadow-utils} محسوب می‌شوند که وظیفه مدیریت متمرکز کاربران و گروه‌ها را در توزیع‌های لینوکس بر عهده دارد. این بسته شامل ابزارهای قدرتمندی است که به مدیران سیستم (\en{Sysadmins}) اجازه می‌دهد نظارت دقیق و همه‌جانبه‌ای بر حساب‌های کاربری داشته باشند.

جدول زیر پرکاربردترین فرامین این مجموعه را معرفی می‌کند:

\begin{table}[htbp]
\centering
\caption{دستورات مجموعه \cmd{shadow-utils} برای مدیریت کاربران و گروه‌ها}
\label{tab:shadow-utils}
\begin{tabularx}{\textwidth}{l X}
\toprule
\textbf{دستور} & \textbf{شرح عملکرد} \\
\midrule
\cmd{lastlog} & گزارش آخرین زمان ورود (\en{Login}) تمامی کاربران یا یک کاربر خاص را نمایش می‌دهد. \\
\addlinespace
\cmd{useradd} & جهت تعریف حساب کاربری جدید یا بازتنظیم پیکربندی‌های پیش‌فرض برای کاربران آتی به کار می‌رود. \\
\addlinespace
\cmd{userdel} & برای حذف کامل یک حساب کاربری و در صورت لزوم، تمامی فایل‌های متعلق به آن استفاده می‌شود. \\
\addlinespace
\cmd{usermod} & جهت اصلاح ویژگی‌های یک حساب موجود (مانند تغییر نام کاربری، دایرکتوری خانگی یا گروه‌ها) کاربرد دارد. \\
\addlinespace
\cmd{groupadd} & وظیفه ایجاد و تعریف یک گروه جدید در لایه سیستم‌عامل را بر عهده دارد. \\
\addlinespace
\cmd{groupdel} & یک گروه موجود را از پایگاه داده گروه‌های سیستم حذف می‌کند. \\
\addlinespace
\cmd{groupmod} & برای ویرایش مشخصات یک گروه (نظیر تغییر نام یا شناسه گروهی \en{GID}) استفاده می‌شود. \\
\bottomrule
\end{tabularx}
\end{table}

این ابزارها ستون فقرات مدیریت هویت و دسترسی در محیط‌های سرور و سیستم‌های چندکاربره (\en{Multi-user}) هستند. برای تسلط بر جزئیات فنی و تمامی گزینه‌های عملیاتی، مطالعه صفحات راهنما از طریق فرمان \cmd{man <command>} توصیه می‌شود.

\section{جمع‌بندی}

در این فصل، معماری سیستم‌های شبه‌یونیکس را از منظر مدیریت سطوح دسترسی (\en{Permissions}) مورد واکاوی قرار دادیم. مدل سنتی مجوزها در لینوکس که بر پایه سه عملیاتِ «خواندن، نوشتن و اجرا» (\cmd{r}، \cmd{w}، \cmd{x}) برای سه سطحِ «مالک، گروه و سایرین» (\cmd{u}، \cmd{g}، \cmd{o}) بنا شده، از نخستین روزهای پیدایش یونیکس تاکنون کارآمدی و پایداری خود را به اثبات رسانده است.

با وجود این، شایان ذکر است که این مکانیزم سنتی در مقایسه با سیستم‌های مدرن‌ترِ مدیریت دسترسی (مانند \en{ACL}ها)، ممکن است در سناریوهای بسیار پیچیده فاقد انعطاف‌پذیری و ظرافت لازم به نظر برسد. با این حال، سادگی ساختاری، ضریب اطمینان بالا و کارایی عملیاتی این مدل سبب شده است که همچنان به عنوان یکی از ارکان بنیادین امنیت در توزیع‌های مختلف لینوکس شناخته شود.

\section{منابع مطالعه تکمیلی}

برای درک عمیق‌تر تهدیدات امنیتی و ماهیت بدافزارها (\en{Malware})، مقاله زیر از دانشنامه ویکی‌پدیا پیشنهاد می‌شود:

\begin{latin}
\url{http://en.wikipedia.org/wiki/Malware}
\end{latin}